% ============================================================
% Relatório técnico (documento vivo) do projeto
% "Da rede lexical à rede contextual" -- MO438 Redes Complexas, Unicamp.
%
% Compilação:  cd report/relatorio && latexmk -pdf relatorio.tex
%
% Material gerado (NÃO editar à mão; ver README.md desta pasta):
%   figuras/*.pdf        -- escritas pela etapa `report` do pipeline
%                           (e, só as de viabilidade, por scripts/feasibility/plot_feasibility.py)
%   tabelas/*.tex        -- escritas pela etapa `report`
%   tabelas/numeros.tex  -- macros com os números citados no texto
% Tudo o que ainda não foi gerado aparece como um quadro "pendente", e o documento
% compila antes de qualquer execução do pipeline.
% ============================================================
\documentclass[11pt,a4paper]{article}

% ------------------------------------------------------------
% Pacotes
% ------------------------------------------------------------
\usepackage[utf8]{inputenc}
\usepackage[T1]{fontenc}
\usepackage[brazil]{babel}
\usepackage{lmodern}
\usepackage{microtype}
\usepackage{amsmath,amssymb}
\usepackage{graphicx}
\usepackage{booktabs}
\usepackage{tabularx}
\usepackage{longtable}
\usepackage{array}
\usepackage[table]{xcolor}
\usepackage{enumitem}
\usepackage{geometry}
\usepackage{titlesec}
\usepackage{caption}
\usepackage{tikz}
\usepackage{seqsplit}  % quebra de linha dentro do SHA da revisão do modelo
\usetikzlibrary{arrows.meta,positioning}
% Sem siunitx de propósito: tabelas/numeros.tex e tabelas/*.tex trazem os números como texto
% simples já formatado em pt-BR (ex.: 15.104; 0,83). Com siunitx, \num{15.104} leria o ponto
% como separador decimal e imprimiria 15,104.
\usepackage{hyperref}

% ------------------------------------------------------------
% Layout (o mesmo da proposta, com seções numeradas)
% ------------------------------------------------------------
\geometry{top=2.0cm, bottom=2.0cm, left=2.0cm, right=2.0cm}

\setlength{\parindent}{0pt}
\setlength{\parskip}{4pt}
% Caminhos e nomes em \texttt não se dividem; um pouco de folga evita linhas estouradas.
\setlength{\emergencystretch}{2.5em}

\setlist[itemize]{leftmargin=1.35em, itemsep=1pt, topsep=2pt, parsep=0pt}
\setlist[enumerate]{leftmargin=1.55em, itemsep=1pt, topsep=2pt, parsep=0pt}
\setlist[description]{leftmargin=1.2em, itemsep=2pt, topsep=2pt, parsep=0pt}

\titleformat{\section}{\bfseries\large}{\thesection}{0.6em}{}
\titlespacing*{\section}{0pt}{12pt}{4pt}
\titleformat{\subsection}{\bfseries\normalsize}{\thesubsection}{0.6em}{}
\titlespacing*{\subsection}{0pt}{8pt}{2pt}
\titleformat{\subsubsection}{\bfseries\itshape\normalsize}{\thesubsubsection}{0.6em}{}
\titlespacing*{\subsubsection}{0pt}{6pt}{1pt}

\captionsetup{font=small, labelfont=bf, margin=0.5cm}
\captionsetup[table]{position=top, skip=4pt}

\hypersetup{
    colorlinks=true,
    linkcolor=black,
    citecolor=black,
    urlcolor=blue,
    pdftitle={Da rede lexical à rede contextual: relatório técnico},
    pdfauthor={Juliana Midlej},
}

\setlength{\LTpre}{4pt}
\setlength{\LTpost}{4pt}

% ------------------------------------------------------------
% Números gerados pela etapa `report`
% ------------------------------------------------------------
% tabelas/numeros.tex define (com \newcommand) as macros dos números citados no texto.
% Os \providecommand abaixo só valem enquanto o arquivo não existe ou não define a macro:
% assim o texto compila antes de qualquer execução e nunca traz um número digitado à mão.
\newcommand{\pendente}{\textit{[pendente]}}
\IfFileExists{tabelas/numeros.tex}{% Gerado pela etapa `report` (gender_networks.report). Não editar à mão.

% Macros dos números citados no texto; relatorio.tex usa \providecommand para as que
% faltarem (impressas como [pendente]).
\newcommand{\nArtigosCandidatos}{2.631}
\newcommand{\nArtigos}{2.153}
\newcommand{\nParagrafos}{17.066}
\newcommand{\nPalavrasCorpus}{1.509.556}
\newcommand{\nTokensCorpus}{2.657.592}
\newcommand{\nTiposCorpus}{24.667}
\newcommand{\nVertices}{13.606}
\newcommand{\nTipos}{2.864}
\newcommand{\nSequencias}{3.230}
\newcommand{\nTokensProcessados}{324.356}
\newcommand{\nParagrafosNucleo}{80}
\newcommand{\nVerticesNucleo}{10.020}
\newcommand{\nVerticesAlvo}{477}
\newcommand{\nVerticesControle}{277}
\newcommand{\nVerticesMultitema}{2.832}
\newcommand{\nAlvosMantidos}{11}
\newcommand{\nGruposPrefixo}{24}
\newcommand{\nTiposVocab}{151.643}
\newcommand{\nTiposVocabCorpus}{24.667}
\newcommand{\difMaxPrefixo}{448,0}
\newcommand{\concordanciaPredicao}{100,0\,\%}
\newcommand{\dataExecucao}{05/10/2026}
\newcommand{\revisaoGit}{dd62bdd-{}dirty (etapas em revisões diferentes; ver Tabela~\ref{tab:versoes})}
}{}
% corpus (etapa corpus; tokens e tipos contados pela etapa sample)
\providecommand{\nArtigosCandidatos}{\pendente}
\providecommand{\nArtigos}{\pendente}
\providecommand{\nParagrafos}{\pendente}
\providecommand{\nPalavrasCorpus}{\pendente}
\providecommand{\nTokensCorpus}{\pendente}
\providecommand{\nTiposCorpus}{\pendente}
% amostra (etapa sample)
\providecommand{\nVertices}{\pendente}
\providecommand{\nTipos}{\pendente}
\providecommand{\nSequencias}{\pendente}
\providecommand{\nTokensProcessados}{\pendente}
\providecommand{\nParagrafosNucleo}{\pendente}
\providecommand{\nVerticesNucleo}{\pendente}
\providecommand{\nVerticesAlvo}{\pendente}
\providecommand{\nVerticesControle}{\pendente}
\providecommand{\nVerticesMultitema}{\pendente}
\providecommand{\nAlvosMantidos}{\pendente}
\providecommand{\nGruposPrefixo}{\pendente}
% vocabulário (etapas sample e knn)
\providecommand{\nTiposVocab}{\pendente}
\providecommand{\nTiposVocabCorpus}{\pendente}
% extração (etapa extract)
\providecommand{\difMaxPrefixo}{\pendente}
\providecommand{\concordanciaPredicao}{\pendente}
% execução
\providecommand{\dataExecucao}{\pendente}
\providecommand{\revisaoGit}{\pendente}

% ------------------------------------------------------------
% Figuras e tabelas opcionais (geradas pelo pipeline)
% ------------------------------------------------------------
\newsavebox{\caixafigura}
% Inclui um PDF no tamanho natural, reduzindo só se for mais largo que a coluna
% (as figuras são desenhadas na largura final, então as fontes não mudam de tamanho).
\newcommand{\incluirlimitado}[1]{%
    \sbox{\caixafigura}{\includegraphics{#1}}%
    \ifdim\wd\caixafigura>\linewidth
        \includegraphics[width=\linewidth]{#1}%
    \else
        \usebox{\caixafigura}%
    \fi}
% Quadro que ocupa o lugar de um material ainda não gerado.
\newcommand{\caixapendente}[1]{%
    \fcolorbox{black!45}{black!4}{%
        \parbox{0.9\linewidth}{\centering\small\color{black!75}%
            \rule{0pt}{1.8em}#1\par\vspace{1.1em}}}}
\newcommand{\etapa}[1]{\texttt{#1}}
% \figuraopcional[quem gera]{arquivo sem extensão}{legenda}{rótulo}
\newcommand{\figuraopcional}[4][etapa \etapa{report}]{%
    \begin{figure}[htbp]
        \centering
        \IfFileExists{figuras/#2.pdf}{%
            \incluirlimitado{figuras/#2.pdf}%
        }{%
            \caixapendente{\textbf{Figura pendente:}
                \texttt{\detokenize{figuras/#2.pdf}}\\[2pt]
                Gerada por: #1.}%
        }
        \caption{#3}\label{#4}
    \end{figure}}
% \tabelaopcional[quem gera]{arquivo sem extensão}{legenda}{rótulo}
% O arquivo gerado contém só o ambiente tabular (booktabs); legenda e rótulo ficam aqui.
\newcommand{\tabelaopcional}[4][etapa \etapa{report}]{%
    \begin{table}[htbp]
        \centering
        \caption{#3}\label{#4}
        \IfFileExists{tabelas/#2.tex}{%
            \small\input{tabelas/#2.tex}%
        }{%
            \caixapendente{\textbf{Tabela pendente:}
                \texttt{\detokenize{tabelas/#2.tex}}\\[2pt]
                Gerada por: #1.}%
        }
    \end{table}}

% ------------------------------------------------------------
% Notação
% ------------------------------------------------------------
\newcommand{\tok}[1]{\texttt{#1}}              % um token, como aparece no texto
\newcommand{\esp}{\textvisiblespace}           % espaço inicial de um token
\newcommand{\cod}[1]{\texttt{#1}}              % código, arquivos, comandos
\newcommand{\eps}{\varepsilon}
\newcommand{\N}{\mathcal{N}}
\newcommand{\previsto}[1]{%
    \par\noindent\fcolorbox{black!30}{white}{\parbox{\dimexpr\linewidth-2\fboxsep-2\fboxrule}{%
        \small\textbf{Conteúdo previsto.} #1}}\par}

% ============================================================
\begin{document}

\begin{center}
    {\large\bfseries MO438 -- Redes Complexas (Unicamp)}\\[4pt]
    {\LARGE\bfseries Da rede lexical à rede contextual}\\[3pt]
    {\large Relatório técnico do experimento com o Qwen3-4B-Base}\\[8pt]
    Juliana Midlej\\[2pt]
    {\small Documento vivo, escrito junto com a implementação. Compilado em \today.
    Números da execução de \dataExecucao{} (\texttt{git} \revisaoGit).}
\end{center}

\begin{abstract}
\noindent
Este relatório registra, etapa por etapa, o experimento que representa a passagem das
representações lexicais às contextuais de um Transformer como uma sequência de redes complexas.
Os vértices são \emph{ocorrências} de tokens em textos da Wikipédia em português; as arestas
ligam representações semelhantes por $k$ vizinhos mais próximos (k-NN) sob o cosseno. Os
mesmos vértices recebem vetores do \emph{embedding} de entrada e das saídas dos blocos 1, 18 e
36 do Qwen3-4B-Base, o que permite comparar as redes vértice a vértice. O texto explica o que
foi feito e por quê, as alternativas descartadas, os problemas encontrados e as limitações.
Figuras, tabelas e números vêm diretamente dos artefatos da última execução do pipeline; o que
ainda não foi executado aparece marcado como \pendente.
\end{abstract}

\tableofcontents

% ============================================================
\section{Introdução e motivação}\label{sec:introducao}
% ============================================================

Transformers representam cada token por um vetor num espaço de alta dimensão. Na entrada do
modelo, esse vetor é uma linha da matriz de \emph{embeddings}: todas as ocorrências de um mesmo
tipo de token começam com \emph{exatamente} o mesmo vetor, não importa o texto em que
apareçam. À medida que atravessam os blocos do Transformer, as representações incorporam
informação do contexto (atenção sobre os tokens anteriores e transformações não lineares) e
se reorganizam no espaço vetorial. No último bloco, a representação de uma ocorrência serve
para prever o \emph{próximo} token, e não mais para identificar o token atual.

A pergunta de fundo deste projeto é \textbf{como} essa reorganização acontece: quanto da
estrutura de similaridade entre representações continua organizada pela identidade lexical e
quanto passa a refletir o contexto, e em que camadas. Estudos anteriores mediram essa
passagem com estatísticas de pares de vetores, como a autossimilaridade das ocorrências de uma
palavra e a anisotropia do espaço~\cite{ethayarajh2019}. Essas medidas resumem distâncias, mas
não dizem \emph{quem é vizinho de quem}, nem se a mudança é só local (troca de alguns vizinhos)
ou se altera a organização em escala maior (\emph{hubs}, comunidades, distâncias).

\subsection{Por que redes complexas}

A proposta é representar cada espaço de representações como uma \textbf{rede}: os vértices são
ocorrências de tokens (ou tipos, na rede lexical de tipos) e cada vértice se liga aos seus $k$
vizinhos mais próximos pelo cosseno. Essa construção tem três vantagens para as perguntas do
projeto:
\begin{itemize}
    \item \textbf{Comparação vértice a vértice.} As redes por ocorrência têm \emph{os mesmos
    vértices} em todas as representações (lexical e camadas 1, 18 e 36). A vizinhança de uma
    ocorrência pode ser comparada diretamente entre camadas (Jaccard), e a fração de vizinhos do
    mesmo tipo (dominância lexical) mede quanto a identidade do token ainda manda.
    \item \textbf{Escalas diferentes.} As ferramentas de redes complexas separam a escala local
    (vizinhança, clusterização) da global (distribuição de graus, \emph{hubs}, distâncias,
    componentes e comunidades). Isso permite distinguir uma reorganização que só troca vizinhos
    de uma que muda a estrutura inteira.
    \item \textbf{Rótulos externos.} As comunidades encontradas podem ser comparadas com
    partições conhecidas (tipo de token, tema, artigo, parágrafo, sentença, próximo token) por
    medidas de concordância, o que transforma ``a rede ficou contextual'' numa afirmação
    mensurável.
\end{itemize}

\subsection{O que o experimento faz}

O experimento usa cerca de $1{,}5\times10^4$ ocorrências de tokens em parágrafos da Wikipédia em
português, escolhidas em oito temas, e o modelo aberto Qwen3-4B-Base~\cite{qwen3} (36 blocos,
dimensão 2560). Para cada ocorrência, extraímos o \emph{embedding} de entrada do seu token e
a saída bruta dos blocos 1, 18 e 36 (e, para a curva camada a camada, dos 36 blocos). Com
esses vetores construímos:
\begin{enumerate}
    \item a \textbf{rede lexical de tipos}, sobre o vocabulário inteiro do modelo (cerca de
    151,6 mil tipos), que descreve a geometria do espaço de entrada;
    \item a \textbf{rede lexical por ocorrência}, em que cada ocorrência recebe o vetor do seu
    tipo, de modo que ocorrências do mesmo tipo são idênticas;
    \item as \textbf{redes contextuais}, com os mesmos vértices representados pelos estados
    internos de cada camada analisada.
\end{enumerate}
As redes são construídas por k-NN com similaridades calculadas em \emph{float64}, com três
tratamentos explícitos de empates (os empates são a regra, e não a exceção, na rede lexical
por ocorrência), nas versões direcionada, por união e mútua. As Seções~\ref{sec:redes}
e~\ref{sec:metodos} descrevem esse método em detalhe.

\subsection{Como ler este relatório}

Este é um \textbf{documento vivo}: ele é atualizado a cada etapa do pipeline e registra o que
foi feito, por que foi feito assim, o que deu errado e o que se encontrou. Ele serve de base
para o relatório final da disciplina, que pode ser uma versão condensada dele. Três convenções
valem para o documento inteiro:
\begin{itemize}
    \item Todo número de resultado citado no texto vem de \cod{tabelas/numeros.tex}, gerado
    pela etapa \etapa{report} a partir dos artefatos da última execução; nada é digitado à
    mão. Enquanto a etapa correspondente não roda, o número aparece como \pendente.
    \item Figuras (\cod{figuras/*.pdf}) e tabelas (\cod{tabelas/*.tex}) também são geradas.
    Quando ainda não existem, um quadro indica o arquivo esperado e a etapa que o produz.
    Números de planejamento (estudo de viabilidade, parâmetros de configuração) são fixos e
    aparecem diretamente no texto, com a fonte indicada.
    \item As Seções~\ref{sec:resultados}--\ref{sec:discussao} ainda são só estrutura: cada
    subseção diz o que vai conter e de qual etapa vêm os dados.
\end{itemize}

\begin{table}[htbp]
\centering
\caption{Estado das seções deste relatório.}\label{tab:estado}
\small
\begin{tabularx}{\linewidth}{@{}l l X@{}}
\toprule
\textbf{Seção} & \textbf{Estado} & \textbf{Depende de} \\
\midrule
1--3. Introdução, perguntas, terminologia & escrita & proposta e plano de execução \\
4. Dados & desenho e viabilidade escritos & números do corpus e da amostra: etapas
\etapa{corpus} e \etapa{sample} \\
5. Modelo e extração & método escrito & diagnósticos: etapa \etapa{extract} \\
6--7. Redes e métodos & método escrito & tempos e estatísticas de empates: etapa \etapa{knn} \\
8. Resultados & estrutura & etapas \etapa{knn}, \etapa{metrics}, \etapa{analyze} e
\etapa{lens} \\
9. Robustez & estrutura & grade de variantes das etapas \etapa{metrics} e \etapa{analyze} \\
10. Discussão & estrutura & resultados \\
11--12. Limitações e problemas & escritas até o planejamento & atualizadas a cada etapa \\
Apêndices & registro de decisões e comandos escritos & tabela de versões: etapa
\etapa{report} \\
\bottomrule
\end{tabularx}
\end{table}

% ============================================================
\section{Perguntas de pesquisa e hipóteses}\label{sec:perguntas}
% ============================================================

O objetivo geral, como na proposta, é caracterizar por meio de redes complexas quanto e de que
forma a estrutura de similaridade entre representações de tokens se transforma entre o espaço
lexical de entrada e as representações contextuais das camadas inicial, intermediária e final
de um Transformer. A hipótese central \emph{não} é que a identidade lexical desapareça, mas que
exista uma reorganização progressiva das vizinhanças e das comunidades, cuja intensidade possa
ser medida. O objetivo se desdobra em quatro perguntas.

\begin{description}
    \item[P1. Quanto e de que forma a organização da rede muda?]
    Quantificar a reorganização \emph{local}, pela similaridade de Jaccard entre as vizinhanças
    de uma mesma ocorrência em duas redes e pela dominância lexical (fração de vizinhos do mesmo
    tipo), e a reorganização \emph{global}, pela distribuição de graus de entrada (surgimento ou
    desaparecimento de \emph{hubs}), pela reciprocidade, pela clusterização, pelas distâncias e
    pelas comunidades.

    \emph{Hipótese H1.} Na rede lexical por ocorrência a dominância é máxima por construção
    (ocorrências do mesmo tipo são idênticas). Ela cai ao longo das camadas, mas continua acima
    do acaso no bloco 36. A mudança não é só local: esperamos \emph{hubs} diferentes em cada
    camada (em espaços de alta dimensão, alguns pontos viram vizinhos de muitos
    outros~\cite{radovanovic2010}) e comunidades menos alinhadas ao tipo.

    \item[P2. As mudanças concentram-se em determinadas camadas?]
    Comparar as transições lexical $\to$ bloco 1, bloco 1 $\to$ 18 e bloco 18 $\to$ 36, e as
    acumuladas, identificando a que concentra a maior queda de Jaccard e de dominância lexical.
    A curva camada a camada (os 36 blocos) localiza a mudança com mais precisão. A resposta é
    dada por faixa de frequência e por categoria de token (palavra inteira ou fragmento).

    \emph{Hipótese H2.} Fragmentos de palavras (início e continuação) precisam do contexto
    imediato para ``virar'' uma palavra; esperamos que a vizinhança deles mude mais cedo (já no
    bloco 1) do que a de palavras inteiras de conteúdo. Tokens frequentes, cujo
    \emph{embedding} é pouco informativo sozinho (palavras funcionais), também devem mudar
    cedo.

    \item[P3. Surgem comunidades associadas a contextos semelhantes?]
    Verificar se as comunidades deixam de ser explicadas pelo tipo de token e passam a ser
    explicadas pelo contexto de origem (tema, artigo, parágrafo, sentença) ou pelo que vem
    depois (o próximo token), com medidas de concordância entre partições (NMI, ARI, pureza)
    corrigidas por uma linha de base de permutação, e por inspeção qualitativa.

    \emph{Hipótese H3.} A concordância com o tipo cai e a concordância com o contexto sobe ao
    longo das camadas. No bloco 36, a concordância com o \emph{próximo} token deve crescer,
    porque é essa a informação que a última camada precisa carregar. A faixa de posição entra
    como controle: se ela explicar as comunidades, o resultado é um artefato posicional, e não
    contexto.

    \item[P4. Ocorrências de um mesmo token separam-se conforme o contexto?]
    Para tipos com muitas ocorrências ($f_t\ge 10$), em especial palavras usadas em sentidos
    diferentes em temas diferentes (as \emph{palavras-alvo}, como \emph{banco}, \emph{campo} e
    \emph{órgão}), verificar se representações inicialmente idênticas passam a ocupar
    comunidades ou regiões diferentes da rede.

    \emph{Hipótese H4.} As palavras-alvo se separam mais, por tema de sentido, do que as
    palavras de controle (\emph{ano}, \emph{grande}, \emph{água}\ldots), que têm frequência e
    distribuição por tema comparáveis mas um sentido estável.
\end{description}

\subsection{Como cada pergunta será respondida}

A Tabela~\ref{tab:perguntas} liga cada pergunta às redes, medidas e seções de resultados. Duas
regras valem para todas as respostas:
\begin{itemize}
    \item \textbf{Nada é interpretado abaixo do nível de ruído.} Na rede lexical por
    ocorrência, a escolha entre vizinhos empatados é arbitrária. Uma mudança lexical $\to$
    contextual só conta como reorganização se for maior que a variação produzida pelo próprio
    desempate (o \emph{piso de ruído}, Seção~\ref{sec:met-ruido}). Da mesma forma, NMI só é
    lido depois de descontada a linha de base de permutação.
    \item \textbf{Resultados sempre estratificados.} Faixa de frequência, categoria de token,
    estrato da amostra e faixa de posição mudam os valores esperados das medidas (por exemplo, a
    dominância lexical máxima depende de $f_t$). Por isso as distribuições são reportadas por
    grupo, e não só como médias globais.
\end{itemize}

\begin{table}[htbp]
\centering
\caption{Perguntas, redes usadas, medidas e onde os resultados aparecem.}\label{tab:perguntas}
\small
\begin{tabularx}{\linewidth}{@{}l >{\raggedright\arraybackslash}p{3.6cm} X l@{}}
\toprule
& \textbf{Redes} & \textbf{Medidas} & \textbf{Seção} \\
\midrule
P1 & por ocorrência (lexical, blocos 1, 18, 36), direcionada e por união &
$J_i$ e $D_i$ por faixa, categoria e estrato; CCDF do grau de entrada, \emph{hubs},
reciprocidade; métricas globais; comunidades & \ref{sec:res-p1} \\
P2 & as mesmas, mais as 36 camadas &
$1-J$ acima do piso de ruído e $\Delta D$ por transição; fração de vértices cuja maior mudança
cai em cada transição; curvas $J(\ell,\ell+1)$ e $D(\ell)$ & \ref{sec:res-p2} \\
P3 & só o núcleo da amostra, por união &
NMI (bruto e menos a linha de base), ARI e pureza contra tipo, tema, artigo, parágrafo,
sentença, próximo token e faixa de posição, camada a camada & \ref{sec:res-p3} \\
P4 & rede inteira; tipos com $f_t\ge10$ &
número efetivo de comunidades $e^{H}$, NMI comunidade$\times$tema dentro do tipo,
autossimilaridade, cosseno intra $-$ entre temas, separação por sentido anotado, redes ego &
\ref{sec:res-p4} \\
Vocabulário & rede lexical de tipos (151,6 mil vértices) &
graus, \emph{hubs}, clusterização, distâncias; comunidades contra a classe de escrita;
vizinhança das palavras-alvo & \ref{sec:res-vocab} \\
\bottomrule
\end{tabularx}
\end{table}

% ============================================================
\section{Terminologia: palavras, tokens, tipos e ocorrências}\label{sec:terminologia}
% ============================================================

A análise mistura três níveis (o texto, o tokenizer e o modelo), e muitas conclusões dependem de
não confundi-los. Esta seção fixa os termos usados no resto do relatório.

\subsection{Definições}

\begin{description}
    \item[Palavra.] Sequência máxima de caracteres de palavra (letras, dígitos e sublinhado,
    a classe \verb|\w| do Unicode) no texto, antes da tokenização. Com essa definição,
    ``tornou-se'' tem duas palavras (\emph{tornou} e \emph{se}) e ``1990'' é uma palavra.
    \item[Token.] Unidade produzida pelo tokenizer do modelo. Pode ser uma palavra inteira, um
    fragmento de palavra, um sinal de pontuação, um dígito ou só espaço. O tokenizer do Qwen3 é
    um BPE em bytes: o espaço antes de uma palavra faz parte do token seguinte (\tok{\esp banco}
    é um token só), e cada dígito é um token separado.
    \item[Tipo de token.] Um elemento do vocabulário, identificado pelo \cod{token\_id} e
    denotado por $t$. O vocabulário do tokenizer tem 151.669 ids, dos quais 26 são tokens
    especiais (\tok{<|endoftext|>}, \tok{<|im\_start|>}, \ldots). O número de ocorrências de
    $t$ na amostra é $f_t$; no corpus inteiro, $f^{\mathrm{corpus}}_t$.
    \item[Ocorrência.] Uma instância posicional de um tipo num texto, denotada por $i$, com tipo
    $t_i$. Cada ocorrência escolhida é um \textbf{vértice} das redes por ocorrência e tem um
    \cod{occurrence\_id}, que é também a linha dela em todos os arquivos de representações.
    \item[Sequência.] O que entra no modelo num \emph{forward pass}: o token de prefixo
    \tok{<|endoftext|>} seguido dos tokens de um parágrafo (Seção~\ref{sec:mod-sequencias}).
    A posição 0 é sempre o prefixo, que nunca vira vértice.
\end{description}

\subsection{Categorias de token}\label{sec:categorias}

Cada ocorrência recebe uma categoria, que define os grupos em que as medidas são reportadas:

\begin{description}
    \item[\cod{whole\_word} (palavra inteira).] O token cobre sozinho uma palavra inteira.
    \item[\cod{word\_start} (início de palavra fragmentada).] Primeiro token de uma palavra
    dividida em vários tokens.
    \item[\cod{continuation} (continuação).] Token que continua uma palavra já começada.
    \item[\cod{punctuation} (pontuação).] Token sem caracteres alfanuméricos, mas com algum
    caractere visível (``\tok{\esp(}'', ``\tok{)}'', ``\tok{,}'').
    \item[\cod{number} (número).] Token cujos caracteres alfanuméricos são todos dígitos. Como
    o Qwen3 divide números dígito a dígito, ``1990'' vira quatro tokens de número, que formam
    uma palavra só.
\end{description}

Tokens formados só por espaço (por exemplo, o espaço antes de um número) \textbf{não viram
vértices}: eles não têm conteúdo próprio e apareceriam em quase toda sequência. Eles
continuam no texto como contexto.

A categoria é calculada pelos \emph{offsets} de caracteres de cada token no parágrafo (o trecho
exato do texto que o token cobre), e \emph{não} pela decodificação isolada do id. Num BPE em
bytes, um caractere acentuado pode ser dividido em dois tokens, e cada pedaço, decodificado
sozinho, vira o caractere de substituição U+FFFD; os \emph{offsets} continuam apontando para o
caractere certo. Cada ocorrência também guarda a palavra a que pertence, a posição do token na
palavra, o número de tokens da palavra e se a palavra é funcional (artigos, preposições,
pronomes, conjunções, verbos auxiliares; lista em \cod{tokens.py}). As interpretações
semânticas se concentram em palavras inteiras de conteúdo, porque um fragmento como
\tok{co} não tem significado próprio.

A Tabela~\ref{tab:exemplos-tokens} mostra três casos reais do tokenizer do Qwen3-4B-Base (revisão
fixada), que guiaram as regras acima.

\begin{table}[htbp]
\centering
\caption{Exemplos de tokenização do Qwen3-4B-Base e as categorias atribuídas. O símbolo
\esp{} marca o espaço que faz parte do token.}\label{tab:exemplos-tokens}
\small
\begin{tabular}{@{}l l r l l@{}}
\toprule
\textbf{Texto} & \textbf{Token} & \textbf{id} & \textbf{Categoria} & \textbf{Palavra} \\
\midrule
``\esp 1990'' & \tok{\esp} & 220 & só espaço (não é vértice) & -- \\
 & \tok{1} & 16 & \cod{number} & 1990 \\
 & \tok{9} & 24 & \cod{number} & 1990 \\
 & \tok{9} & 24 & \cod{number} & 1990 \\
 & \tok{0} & 15 & \cod{number} & 1990 \\
\midrule
``\esp tornou-se'' & \tok{\esp torn} & 21145 & \cod{word\_start} & tornou \\
 & \tok{ou} & 283 & \cod{continuation} & tornou \\
 & \tok{-se} & 7806 & \cod{whole\_word} & se \\
\midrule
``\esp(banco)'' & \tok{\esp(} & 320 & \cod{punctuation} & -- \\
 & \tok{ban} & 6850 & \cod{word\_start} & banco \\
 & \tok{co} & 1015 & \cod{continuation} & banco \\
 & \tok{)} & 8 & \cod{punctuation} & -- \\
\bottomrule
\end{tabular}
\end{table}

O terceiro caso mostra por que as palavras-alvo são contadas \textbf{só} no id da forma
minúscula precedida de espaço: \tok{\esp banco} é o id 52465, mas ``(banco)'' vira
\tok{ban}${}+{}$\tok{co}, ``Banco'' é outro tipo (\tok{\esp Banco}, id 76687), e o mesmo vale
para ``Estado'' (id 40174) contra \tok{\esp estado} (id 24062). Considerar essas variantes
misturaria tipos diferentes numa mesma ``palavra'', o que quebraria a premissa da P4 de que
as ocorrências comparadas partem do mesmo vetor lexical.

\subsection{Faixas de frequência e de posição}\label{sec:faixas}

Os tipos são agrupados em faixas de frequência na amostra, $f_t\in\{1\text{--}2\}$,
$\{3\text{--}9\}$, $\{10\text{--}49\}$ e $\{50\}$, onde $50=f_{\max}$ é o limite de ocorrências
por tipo. Como a amostra é em parte desenhada (as palavras-alvo são levadas até perto de 50
ocorrências, as multitema até 20), a frequência na amostra reflete o desenho. Por isso cada
ocorrência também tem a faixa de frequência \emph{no corpus}, com as mesmas fronteiras e o topo
aberto ($50+$), que não depende da amostragem.

Tipos com $f_t\le 2$ continuam como vértices (para que o conjunto de vértices seja o mesmo em
todas as redes), mas ficam fora das análises por tipo: para $f_t=1$ a dominância lexical é nula
por definição, e para $f_t=2$ ela não passa de $1/k$. Dominância e análises intra-tipo usam
$f_t\ge3$, e a P4 usa $f_t\ge10$.

A posição na sequência também é agrupada, em $\{1\text{--}4\}$, $\{5\text{--}16\}$,
$\{17\text{--}64\}$ e $\{65+\}$. As primeiras posições têm estados atípicos (pouco contexto,
normas grandes), e a faixa de posição serve de controle nas análises.

\subsection{Os três tipos de rede e a notação}\label{sec:tipos-rede}

Para uma ocorrência $i$, $e_{t_i}\in\mathbb{R}^{2560}$ é o \emph{embedding} de entrada do seu
tipo e $h_i^{(\ell)}$ é a saída do bloco $\ell$ na posição de $i$. O vetor $x_i$ usado para
construir uma rede é um deles, e $\N_i$ é o conjunto de vizinhos de $i$ nessa rede.
\begin{enumerate}
    \item \textbf{Rede lexical de tipos (vocabulário):} um vértice para cada tipo do
    vocabulário (menos os especiais), representado por $e_t$, com a frequência no corpus e na
    amostra e a classe de escrita como atributos.
    \item \textbf{Rede lexical por ocorrência:} os vértices são as ocorrências da amostra, e
    $x_i=e_{t_i}$. Ocorrências do mesmo tipo têm vetores idênticos.
    \item \textbf{Redes contextuais:} os mesmos vértices, com $x_i=h_i^{(\ell)}$ para
    $\ell\in\{1,18,36\}$ (e as variantes da Seção~\ref{sec:red-robustez}).
\end{enumerate}
Na variante (c) do tratamento de empates (Seção~\ref{sec:red-empates}), a vizinhança de $i$ é
um conjunto $T_i$ de \emph{tipos} distintos, e não de ocorrências.

\subsection{Rede do vocabulário e tabela de tipos da amostra}\label{sec:vocab-vs-amostra}

Há dois objetos lexicais ``de tipos'' no experimento, e eles não são a mesma coisa:
\begin{itemize}
    \item a \textbf{rede lexical de tipos}, com k-NN entre os $\approx$151,6 mil tipos do
    vocabulário. É uma rede reportada nos resultados, com as métricas da disciplina;
    \item a \textbf{tabela de vizinhos dos tipos da amostra}, com k-NN entre os $\approx$3,3 mil
    tipos distintos que aparecem entre as ocorrências da amostra. Ela não é reportada como
    rede. Serve para (i) derivar de forma exata a rede lexical por ocorrência
    (Seção~\ref{sec:red-derivacao}) e (ii) dar o $T_i$ lexical da variante (c).
\end{itemize}

Seria tentador obter a tabela da amostra recortando a rede do vocabulário (o subgrafo induzido
pelos tipos da amostra). Isso não funciona, e a Figura~\ref{fig:recorte} mostra por quê com um
vocabulário de cinco tipos, $A$--$E$, dos quais só $A$, $B$ e $C$ aparecem na amostra, e $k=1$.
As similaridades são
$s(A,D)=0{,}9$, $s(B,E)=0{,}8$, $s(C,D)=0{,}7$, $s(A,B)=0{,}6$, $s(A,E)=0{,}4$,
$s(A,C)=0{,}3$, $s(B,C)=0{,}2$ e as demais menores.
\begin{itemize}
    \item Na rede do vocabulário, o vizinho mais próximo de $A$ é $D$, o de $B$ é $E$ e o de
    $C$ é $D$. Recortando os tipos da amostra, \emph{nenhuma} aresta sobra: todos os vizinhos
    estão fora da amostra.
    \item O k-NN calculado só entre $A$, $B$ e $C$ dá $A\to B$, $B\to A$ e $C\to A$, arestas que
    \emph{não existem} na rede do vocabulário.
\end{itemize}
Logo, nenhuma das duas é subgrafo da outra: o recorte perde os vizinhos que estão fora da
amostra, e não tem as arestas que o k-NN restrito à amostra tem. Para as redes por ocorrência,
a vizinhança certa é a restrita à amostra, porque nas camadas contextuais só existem vetores
para tokens que aparecem no texto; para comparar a camada lexical com as contextuais de forma
justa, o $T_i$ lexical também precisa ser calculado entre os tipos da amostra.

Isso corrige uma frase da Seção~7.3 da proposta, escrita quando a rede de tipos ainda seria
construída só com os tipos da amostra: \emph{na etapa lexical, $T_i$ coincide com o k-NN entre
os tipos da amostra} (a tabela auxiliar), e não com a rede lexical de tipos, que agora é a do
vocabulário inteiro.

\begin{figure}[htbp]
\centering
\begin{tikzpicture}[
    every node/.style={font=\small},
    tipo/.style={circle, draw, minimum size=6.5mm, inner sep=0pt},
    amostra/.style={tipo, fill=black!12, very thick},
    fora/.style={tipo, dashed, draw=black!55, text=black!60},
    seta/.style={-{Stealth[length=2.2mm]}, thick},
    titulo/.style={font=\small\bfseries, align=center},
]
% (a) vocabulário
\begin{scope}
    \node[amostra] (A) at (0,0) {$A$};
    \node[amostra] (B) at (1.6,0) {$B$};
    \node[amostra] (C) at (0,-1.6) {$C$};
    \node[fora] (D) at (-1.3,-0.8) {$D$};
    \node[fora] (E) at (2.6,-0.9) {$E$};
    \draw[seta] (A) -- (D);
    \draw[seta] (C) -- (D);
    \draw[seta] (B) -- (E);
    \draw[seta, draw=black!55] (D) to[bend left=30] (A);
    \draw[seta, draw=black!55] (E) to[bend left=30] (B);
    \node[titulo] at (0.6,1.0) {(a) vocabulário, $k=1$};
\end{scope}
% (b) recorte
\begin{scope}[xshift=5.4cm]
    \node[amostra] (A2) at (0,0) {$A$};
    \node[amostra] (B2) at (1.6,0) {$B$};
    \node[amostra] (C2) at (0,-1.6) {$C$};
    \node[titulo] at (0.8,1.0) {(b) recorte de (a)\\[-1pt]nos tipos da amostra};
    \node[font=\footnotesize, text=black!60, align=center] at (0.9,-2.35) {nenhuma aresta};
\end{scope}
% (c) k-NN restrito
\begin{scope}[xshift=9.6cm]
    \node[amostra] (A3) at (0,0) {$A$};
    \node[amostra] (B3) at (1.6,0) {$B$};
    \node[amostra] (C3) at (0,-1.6) {$C$};
    \draw[seta] (A3) to[bend left=18] (B3);
    \draw[seta] (B3) to[bend left=18] (A3);
    \draw[seta] (C3) -- (A3);
    \node[titulo] at (0.8,1.0) {(c) k-NN calculado\\[-1pt]só na amostra, $k=1$};
\end{scope}
\end{tikzpicture}
\caption{Por que a tabela de vizinhos dos tipos da amostra não é um recorte da rede do
vocabulário. Tipos cinza estão na amostra; $D$ e $E$ (tracejados) só existem no vocabulário.
O recorte (b) perde todas as arestas de (a), e o k-NN restrito (c) tem arestas que (a) não
tem.}\label{fig:recorte}
\end{figure}

% ============================================================
\section{Dados: corpus, estudo de viabilidade e amostra}\label{sec:dados}
% ============================================================

Os vértices das redes são ocorrências de tokens em parágrafos reais da Wikipédia em português.
Esta seção descreve como o corpus é montado (etapa \etapa{corpus}), o estudo de viabilidade que
dimensionou a amostra e o desenho final da amostra (etapa \etapa{sample}).

% ------------------------------------------------------------
\subsection{Corpus}\label{sec:corpus}

\paragraph{Por que a Wikipédia.} É texto real, em português, com licença aberta e
organizado por categorias, o que permite definir \emph{temas} de forma reproduzível. Cada
artigo tem um identificador de revisão (\cod{revid}); guardando esses identificadores, o mesmo
texto pode ser baixado de novo no futuro, mesmo que o artigo tenha mudado: o comando
\cod{corpus-rebuild} baixa por \cod{revid} as revisões dos artigos mantidos e as limpa com o
mesmo código e na mesma ordem da etapa \etapa{corpus}.

\paragraph{Temas.} São oito, cada um definido pela árvore de uma categoria raiz: Física,
Biologia, Economia, Política, Ciência da computação, Música, Futebol (tema ``esporte'') e
Geografia. O tema cumpre dois papéis: é um dos rótulos de contexto da P3 e serve como
\emph{indicador aproximado do sentido} das palavras-alvo na P4 (\emph{banco} num artigo de
computação tende a ser banco de dados; em economia, instituição financeira). Os temas foram
escolhidos distantes entre si, para que o vocabulário de cada um seja característico, e de
modo que cada palavra-alvo tenha sentidos diferentes em pelo menos dois deles.

\paragraph{Busca por categorias.} A partir de cada raiz, uma busca em largura percorre as
subcategorias até a profundidade 2 e junta até 600 títulos por tema. Só entram páginas do
domínio principal (artigos); subcategorias de manutenção ou de usuário (títulos que contêm
``Artigos'', ``Esboço'', ``Páginas'', ``Predefinições'', ``Wikipédia'', ``Usuário'', ``Listas''
ou ``Categoria:!'') não são seguidas, e títulos ``Lista de\ldots'' são descartados
(Seção~\ref{sec:prob-categorias}). Todas as respostas cruas da API ficam num cache em disco
(\cod{data/raw/wiki}), o que torna o download retomável e a limpeza refazível sem rede. No
modo offline (\cod{GENDER\_NETWORKS\_OFFLINE=1}), uma resposta que falta no cache interrompe a
etapa, em vez de baixar uma revisão mais nova.

\paragraph{Categorias extras para os sentidos fracos.} No estudo de viabilidade, alguns
sentidos das palavras-alvo apareciam muito pouco no tema que os representa: \emph{nota} em
Economia (7 ocorrências, no sentido de cédula), \emph{carga} em Economia (6, carga
tributária), \emph{órgão} em Biologia (6) e \emph{banco} em Geografia (4, banco de areia);
\emph{matriz} e \emph{planta} também tinham um segundo sentido fraco. Para esses casos, o
corpus inclui categorias extras ligadas a um dos oito temas (por exemplo, Numismática e Impostos
em Economia; Anatomia humana e Botânica em Biologia; Geomorfologia e Oceanografia em
Geografia; Bancos de dados e Matrizes em Computação), com busca até a profundidade 1 e até 60
artigos por categoria. A lista completa está em \cod{configs/experiment.yaml}.

\paragraph{Páginas descartadas.} Um artigo é descartado quando:
\begin{itemize}
    \item é alcançado a partir de \textbf{mais de um tema} (contando as categorias extras). A
    busca registra todos os temas que chegam a cada título. Uma página multitema (por exemplo,
    Economia política) misturaria o vocabulário de dois temas e borraria justamente o rótulo
    usado na P3 e na P4;
    \item é uma \textbf{biografia}, detectada pelas predefinições de infobox
    (\cod{Info/Biografia}, \cod{Info/Futebolista}, \cod{Info/Político}, \cod{Info/Músico} e
    outras). Biografias aparecem em quase todas as árvores de categorias e falam de datas,
    lugares e carreiras, um vocabulário que não é do tema;
    \item é uma página de \textbf{desambiguação}, não tem texto, repete um artigo já escolhido
    ou não sobra nenhum parágrafo depois da limpeza.
\end{itemize}
Os artigos restantes de cada tema são embaralhados com semente fixa e baixados em lotes de 50
até completar a cota de 150 por tema (60 por categoria extra).

\paragraph{Limpeza.} As seções de material de apoio do \emph{wikitext} (Referências, Ligações
externas, Ver também, Bibliografia, Notas, Fontes e variantes) e suas subseções são removidas;
seções de conteúdo que vêm depois delas são mantidas (em artigos de música, por exemplo,
``Notas'' pode vir antes de ``História''). O texto é então convertido em texto com o \cod{mwparserfromhell}. Antes da conversão, a árvore sintática é
reescrita em três pontos, porque a conversão padrão (\cod{strip\_code}) mantém o conteúdo de
marcas que não são texto corrido:
\begin{itemize}
    \item marcas de bloco (\cod{<ref>}, \cod{<gallery>}, \cod{<timeline>}, \cod{<pre>},
    tabelas e outras) e links para arquivos e categorias são removidos sem deixar rastro; sem
    isso, o texto das referências entrava nos parágrafos (Seção~\ref{sec:prob-ref});
    \item marcas que não podem virar palavras (\cod{<math>}, \cod{<chem>}, \cod{<score>}) e itens
    de lista fazem a linha inteira ser descartada: apagar uma fórmula deixaria um buraco na
    frase (``a energia é dada por , onde\ldots''); fórmulas triviais (um símbolo ou um número,
    como \cod{<math>v</math>}) são mantidas como texto;
    \item predefinições em linha cujo texto visível é um argumento (\cod{formatnum},
    \cod{nowrap}, conversões de unidades como \cod{converter}) são trocadas por esse texto; as
    demais predefinições (infoboxes, navegação, pronúncia, datas formatadas) são descartadas, e
    os restos que deixam (``( )'', ``, ,'') são limpos por expressões regulares.
\end{itemize}

\paragraph{Filtros de parágrafo.} Fica só a prosa corrida: parágrafos com pelo menos 40
palavras, que terminam como frase (``.'', ``!'' ou ``?'', eventualmente seguidos de aspas ou
parênteses), não começam com marcação de tabela ou lista e têm no máximo 8\% de dígitos
(o que elimina tabelas de resultados e cronologias). Parágrafos repetidos são removidos pelos
primeiros 80 caracteres em minúsculas (textos copiados entre artigos são comuns). Cada
parágrafo é dividido em sentenças por expressão regular, com uma lista de abreviações do
português (``Sr.'', ``séc.'', ``p.ex.'', iniciais de nomes\ldots), para dar o rótulo de
sentença usado na P3.

\paragraph{Saídas.} \cod{data/\allowbreak corpus/\allowbreak articles.jsonl} (todos os candidatos, com o motivo de
descarte), \cod{data/\allowbreak corpus/\allowbreak paragraphs.jsonl} (parágrafos mantidos, com os limites das
sentenças) e \cod{data/\allowbreak corpus/\allowbreak manifest.csv}, a lista versionada de \cod{pageid},
\cod{revid}, título, tema e categoria de origem.

\subsubsection*{Estatísticas do corpus}

O corpus final tem \nArtigos{} artigos (de \nArtigosCandidatos{} candidatos), \nParagrafos{}
parágrafos, \nPalavrasCorpus{} palavras, \nTokensCorpus{} tokens do Qwen3 e \nTiposCorpus{}
tipos distintos.

\tabelaopcional[etapa \etapa{report} a partir de \etapa{corpus}]{corpus_temas}
    {Artigos, parágrafos e palavras por tema, e artigos descartados por motivo.}
    {tab:corpus-temas}
\figuraopcional[etapa \etapa{report} a partir de \etapa{corpus}]{corpus_paragrafos}
    {Parágrafos por tema e distribuição do tamanho dos parágrafos (em palavras).}
    {fig:corpus-paragrafos}
\figuraopcional[etapa \etapa{report} a partir de \etapa{corpus}]{corpus_descartes}
    {Páginas candidatas descartadas por motivo (multitema, biografia, desambiguação, sem
    parágrafos, duplicada), por tema.}
    {fig:corpus-descartes}

% ------------------------------------------------------------
\subsection{Estudo de viabilidade}\label{sec:viabilidade}

Antes de implementar o pipeline, um estudo rápido (scripts em \cod{scripts/feasibility/})
respondeu a três perguntas: (i) a meta de cerca de 15 mil vértices é alcançável com parágrafos
inteiros? (ii) Parágrafos inteiros deixam palavras repetidas suficientes, em temas diferentes,
para a P4? (iii) Que tokenizer, e portanto que modelo, usar?

\paragraph{Dados do estudo.} O \cod{fetch\_wiki.py} baixou uma versão simplificada do corpus:
os mesmos oito temas, busca até a profundidade 2, 120 artigos por tema, limpeza só com
\cod{strip\_code} e os mesmos filtros de parágrafo. Resultaram \textbf{958 artigos} e
\textbf{7.876 parágrafos} (Biologia 965, Computação 1.034, Economia 1.147, Esporte 836, Física
948, Geografia 785, Música 1.032, Política 1.129), com \textbf{1,28 milhão de tokens} do
Qwen3 (1.276.475) e cerca de 23 mil tipos distintos. Três tokenizers foram comparados: o do
Qwen3-4B (idêntico ao do Qwen3-4B-Base, 151.669 ids), o do Tucano-2b4 (treinado em português) e
o do GPT-2 small português (o modelo do piloto).

\paragraph{Amostragem natural.} O \cod{simulate\_sample.py} simula a amostragem mais simples:
parágrafos inteiros, em rodízio entre os temas, com todos os tokens como vértices e um limite
$f_{\max}$ por tipo, até chegar a 15 mil vértices. Também simula uma variante
``direcionada'', que começa pelos parágrafos que contêm as palavras-alvo (até 12 por palavra e
tema) e depois segue a ordem natural. As saídas estão em \cod{sim\_<modelo>.txt}.

\paragraph{Desenho híbrido.} O \cod{sim\_hybrid.py} simula a alternativa adotada: um
\emph{núcleo} denso de parágrafos inteiros mais estratos esparsos de ocorrências escolhidas
(palavras-alvo, controles e palavras de conteúdo presentes em vários temas), em que só as
posições sorteadas viram vértices e o resto do parágrafo é contexto. As saídas estão em
\cod{hyb\_<modelo>.txt}.

Nas contagens a seguir, \emph{palavra de conteúdo} é um tipo cuja categoria mais comum é
palavra inteira, só com letras, com pelo menos 3 letras e fora de uma lista de palavras
funcionais. O critério da P4 é ter $f_t\ge10$ na amostra e pelo menos 3 ocorrências em pelo
menos 3 temas.

\begin{table}[htbp]
\centering
\caption{Estudo de viabilidade: tokenizers e amostragem natural até 15 mil vértices com
$f_{\max}=50$. Fontes: linha \cod{pool:} e bloco \cod{f\_max=50, amostragem natural} de
\cod{scripts/feasibility/sim\_<modelo>.txt}.}\label{tab:viab-natural}
\small
\begin{tabular}{@{}l r r r@{}}
\toprule
& \textbf{Qwen3-4B} & \textbf{Tucano-2b4} & \textbf{GPT-2 small PT} \\
\midrule
Tokens no \emph{pool} de 7.876 parágrafos & 1.276.475 & 953.290 & 939.895 \\
Tokens por palavra & 1,78 & 1,33 & 1,31 \\
Palavras-alvo candidatas que são um token só & 12 de 16 & 16 de 16 & 16 de 16 \\
\midrule
Parágrafos para chegar a 15 mil vértices & 101 & 134 & 160 \\
Tipos distintos na amostra & 3.693 & 6.144 & 6.213 \\
Vértices com $f_t\le2$ & 19\% & 40\% & 39\% \\
Vértices que são palavras inteiras & 28\% & 54\% & 72\% \\
Tipos com $f_t\ge10$ & 293 & 183 & 170 \\
Conteúdo com $f_t\ge10$ e $\ge3$ oc.\ em $\ge3$ temas & 4 & 5 & 9 \\
Alvos com $\ge10$ oc.\ em $\ge2$ temas & 0 & 0 & 0 \\
\midrule
\multicolumn{4}{@{}l}{\emph{Variante direcionada (alvos primeiro), $f_{\max}=50$}} \\
Conteúdo com $f_t\ge10$ e $\ge3$ oc.\ em $\ge3$ temas & 10 & 10 & 21 \\
Alvos com $\ge10$ oc.\ em $\ge2$ temas & 2 & 1 & 2 \\
\bottomrule
\end{tabular}
\end{table}

\figuraopcional[\cod{scripts/feasibility/plot\_feasibility.py}]{viabilidade_tokenizers}
    {Comparação dos tokenizers no \emph{pool} do estudo de viabilidade: tokens por palavra e
    fração dos vértices que são palavras inteiras na amostragem natural ($f_{\max}=50$). O
    Qwen3 (em destaque) é o tokenizer do modelo escolhido. Fonte:
    \cod{scripts/feasibility/sim\_<modelo>.txt}.}
    {fig:viab-tokenizers}

A Tabela~\ref{tab:viab-natural} e a Figura~\ref{fig:viab-tokenizers} mostram três coisas.

\textbf{A meta de vértices é fácil.} Com qualquer tokenizer, 15 mil vértices saem de 100 a 160
parágrafos (de 18 a 21 mil tokens), uma fração pequena do \emph{pool}. Com $f_{\max}=20$ bastam
de 127 a 177 parágrafos.

\textbf{Parágrafos inteiros não servem para a P4.} Mesmo com 15 mil vértices, só de 4 a 9
palavras de conteúdo têm $f_t\ge10$ com ocorrências espalhadas por 3 temas, e \emph{nenhuma}
palavra-alvo chega a 10 ocorrências em 2 temas. O motivo é a lei de Zipf: num trecho pequeno de
texto, só palavras funcionais e pontuação se repetem muito; palavras de conteúdo aparecem uma
ou duas vezes. Começar pelos parágrafos que contêm as palavras-alvo (variante direcionada)
melhora pouco (de 10 a 21 palavras de conteúdo, 1 ou 2 alvos).

\textbf{O Qwen3 fragmenta muito o português.} São 1,78 tokens por palavra, contra 1,33 do
Tucano e 1,31 do GPT-2 português, e só 28\% dos vértices são palavras inteiras (54\% e 72\% nos
outros). Quatro das 16 palavras-alvo candidatas não são um token só no Qwen3 (\emph{corrente},
\emph{onda}, \emph{célula} e \emph{núcleo}; por exemplo, ``\esp célula'' vira
\tok{\esp cél}${}+{}$\tok{ula}) e saíram da lista, que ficou com 12. Essa foi a principal
desvantagem pesada contra o Qwen3 na escolha do modelo (Seção~\ref{sec:mod-escolha}).

\begin{table}[htbp]
\centering
\caption{Estudo de viabilidade: desenho híbrido com 15 mil vértices e $f_{\max}=50$
(núcleo de parágrafos inteiros, 12 ou 16 alvos, 6 controles e 150 palavras multitema com 20
ocorrências cada). Fonte: \cod{scripts/feasibility/hyb\_<modelo>.txt}.}\label{tab:viab-hibrido}
\small
\begin{tabular}{@{}l r r r@{}}
\toprule
& \textbf{Qwen3-4B} & \textbf{Tucano-2b4} & \textbf{GPT-2 small PT} \\
\midrule
Vértices & 15.104 & 15.002 & 15.099 \\
\quad núcleo & 11.251 & 10.954 & 11.045 \\
\quad alvos & 553 & 748 & 754 \\
\quad controles & 300 & 300 & 300 \\
\quad multitema & 3.000 & 3.000 & 3.000 \\
Parágrafos inteiros no núcleo & 78 & 93 & 107 \\
Vértices por parágrafo do núcleo (média) & 144 & 118 & 103 \\
Parágrafos tocados no total & 2.894 & 2.943 & 3.017 \\
Tokens a processar no modelo & 329.555 & 243.959 & 248.579 \\
Tipos distintos & 3.257 & 4.904 & 4.856 \\
Vértices com $f_t\le2$ & 17\% & 32\% & 30\% \\
Tipos com $f_t\ge10$ & 354 & 274 & 286 \\
Conteúdo com $f_t\ge10$ e $\ge3$ oc.\ em $\ge3$ temas & 169 & 172 & 173 \\
Alvos com $\ge10$ oc.\ em $\ge2$ temas & 3 de 12 & 7 de 16 & 7 de 16 \\
Alvos com $\ge5$ oc.\ em $\ge3$ temas & 11 de 12 & 14 de 16 & 14 de 16 \\
\bottomrule
\end{tabular}
\end{table}

\figuraopcional[\cod{scripts/feasibility/plot\_feasibility.py}]{viabilidade_amostragem}
    {Amostragem natural (parágrafos inteiros, $f_{\max}=50$) contra o desenho híbrido, nas duas
    contagens que decidem a P4: palavras de conteúdo repetidas em vários temas e palavras-alvo
    com ocorrências suficientes em dois temas. Fontes:
    \cod{scripts/feasibility/sim\_<modelo>.txt} e \cod{hyb\_<modelo>.txt}.}
    {fig:viab-amostragem}

\textbf{O desenho híbrido resolve.} Com o mesmo número de vértices, o núcleo continua com cerca
de 11 mil vértices em 78 a 107 parágrafos inteiros (o bastante para a P3 e para as métricas
globais), e o número de palavras de conteúdo com $f_t\ge10$ espalhadas por 3 temas sobe de
4--9 para cerca de 170 (Tabela~\ref{tab:viab-hibrido} e Figura~\ref{fig:viab-amostragem}). O
custo é processar mais texto no modelo: cerca de 330 mil tokens com o Qwen3, porque cada
ocorrência esparsa exige o parágrafo até a posição dela.

Na simulação, os alvos ainda ficaram fracos (3 de 12 com $\ge10$ ocorrências em 2 temas), mas
por uma razão corrigível: as 50 ocorrências de cada alvo foram sorteadas em rodízio entre os
\emph{oito} temas, e não concentradas nos temas em que a palavra tem sentidos diferentes. A
Tabela~\ref{tab:viab-pool} mostra que o \emph{pool} tem ocorrências suficientes para cotas por
\emph{tema de sentido} na maioria dos casos, e onde faltam (\emph{banco} em Geografia,
\emph{órgão} em Biologia, \emph{nota} e \emph{carga} em Economia, \emph{planta} em Economia,
\emph{matriz} em Computação), as categorias extras da Seção~\ref{sec:corpus} entram. A
simulação também tinha dois defeitos de desenho, corrigidos no desenho final
(Seção~\ref{sec:prob-simulacao}): o limite por tipo do núcleo era aplicado por ordem de
chegada, e cada parágrafo do núcleo vinha de um artigo diferente.

\begin{table}[htbp]
\centering
\caption{Ocorrências das palavras-alvo e de controle no \emph{pool} inteiro do estudo de
viabilidade (tokenizer do Qwen3, sem limite por tipo), por tema. Em negrito, os temas de
sentido usados nas cotas do desenho final (Tabela~\ref{tab:cotas}). Fonte: bloco \cod{pool
inteiro} de \cod{scripts/feasibility/sim\_Qwen3-4B.txt}.}\label{tab:viab-pool}
\footnotesize
\setlength{\tabcolsep}{4.5pt}
\begin{tabular}{@{}l r r r r r r r r r@{}}
\toprule
& \textbf{Total} & \textbf{Bio} & \textbf{Comp} & \textbf{Econ} & \textbf{Esp} & \textbf{Fís}
& \textbf{Geo} & \textbf{Mús} & \textbf{Pol} \\
\midrule
banco & 66 & 9 & \textbf{19} & \textbf{24} & 6 & 0 & \textbf{4} & 3 & 1 \\
campo & 335 & 34 & 26 & 22 & \textbf{74} & \textbf{82} & \textbf{72} & 11 & 14 \\
órgão & 57 & \textbf{6} & 1 & 6 & 9 & 0 & 0 & \textbf{25} & \textbf{10} \\
estado & 282 & 24 & 39 & 28 & 13 & \textbf{84} & 19 & 20 & \textbf{55} \\
carga & 53 & 9 & 4 & \textbf{6} & 0 & \textbf{22} & 0 & 3 & 9 \\
rede & 133 & 4 & \textbf{71} & 15 & 2 & 11 & \textbf{15} & 3 & 12 \\
nota & 41 & 2 & 5 & \textbf{7} & 0 & 2 & 0 & \textbf{24} & 1 \\
capital & 271 & 1 & 0 & \textbf{201} & 18 & 1 & 9 & 6 & \textbf{35} \\
família & 70 & \textbf{10} & 10 & 4 & 4 & 13 & 2 & 13 & \textbf{14} \\
massa & 140 & 2 & 3 & 10 & 3 & \textbf{59} & \textbf{39} & 14 & 10 \\
matriz & 27 & \textbf{10} & \textbf{3} & 2 & 0 & 5 & 0 & 6 & 1 \\
planta & 35 & \textbf{25} & 0 & \textbf{0} & 0 & 0 & 9 & 0 & 1 \\
\midrule
ano & 310 & 13 & 15 & \textbf{36} & \textbf{82} & 35 & 30 & \textbf{65} & 34 \\
século & 356 & 29 & 12 & 31 & 11 & \textbf{86} & \textbf{66} & 59 & \textbf{62} \\
grande & 527 & \textbf{58} & 48 & 67 & 66 & \textbf{88} & 49 & 77 & 74 \\
primeiro & 510 & 24 & \textbf{52} & 39 & \textbf{134} & 63 & 48 & 92 & 58 \\
importante & 247 & \textbf{53} & 20 & \textbf{47} & 9 & 24 & 23 & 36 & 35 \\
água & 249 & \textbf{102} & 0 & 24 & 2 & 47 & \textbf{65} & 1 & 8 \\
\bottomrule
\end{tabular}
\end{table}

Os números do estudo são aproximados por construção: o corpus do estudo não removia páginas
multitema nem biografias, a limpeza ainda deixava restos de referências e as contagens
ignoram os limites por parágrafo e por artigo do desenho final. Eles serviram para escolher o
desenho e o modelo; a composição real da amostra está na Seção~\ref{sec:amostra}.

% ------------------------------------------------------------
\subsection{Amostra híbrida}\label{sec:amostra}

A amostra final tem cerca de 15 mil vértices, $f_{\max}=50$ e sementes fixas (semente 438),
com quatro estratos sorteados em ordem por um único gerador aleatório.

\paragraph{1. Núcleo (P1--P3).} Por tema, 5 artigos $\times$ 2 parágrafos (cerca de 80
parágrafos e 11 mil vértices), com todos os tokens que não são só espaço como vértices. Dois
parágrafos por artigo fazem \emph{artigo} e \emph{parágrafo} serem rótulos distintos na P3 (na
simulação eles quase coincidiam). O limite $f_{\max}$ é aplicado por \textbf{sorteio uniforme}
entre as ocorrências de cada tipo acima do limite, depois de juntar todo o núcleo, e não por
ordem de chegada: assim as ocorrências mantidas de ``de'' ou ``,'' se espalham por todos os
parágrafos, em vez de serem as dos primeiros parágrafos. O núcleo é o único estrato em que um
parágrafo aparece inteiro, e por isso a P3 é analisada só nele.

\paragraph{2. Palavras-alvo (P4).} Doze palavras com sentidos diferentes em temas diferentes:
\emph{banco}, \emph{campo}, \emph{órgão}, \emph{estado}, \emph{carga}, \emph{rede},
\emph{nota}, \emph{capital}, \emph{família}, \emph{massa}, \emph{matriz} e \emph{planta}. Cada
uma tem cotas por \textbf{tema de sentido}: 2 ou 3 temas, com 16 a 25 ocorrências em cada
(Tabela~\ref{tab:cotas}). A cota conta as ocorrências do alvo que já estão no núcleo daquele
tema; só a diferença é sorteada em parágrafos fora do núcleo, com no máximo 1 ocorrência por
(tipo, parágrafo) e 3 por (tipo, artigo), para que um único artigo não domine a amostra de uma
palavra. As ocorrências que ficam no núcleo mantêm o estrato \cod{core}, mas recebem
\cod{target\_word} e \cod{sense\_theme} como as sorteadas; por isso as contagens por estrato
(\nVerticesAlvo{} alvos, \nVerticesControle{} controles) ficam abaixo da soma das cotas, sem que
falte ocorrência. Só conta o id da forma minúscula precedida de espaço
(Seção~\ref{sec:categorias}). Um alvo que não chegar a 10 ocorrências (\cod{min\_per\_sense})
em pelo menos 2 temas sai da lista, e as faltas por (alvo, tema) ficam registradas.

\paragraph{3. Controles (P4).} Seis palavras de sentido estável (\emph{ano}, \emph{século},
\emph{grande}, \emph{primeiro}, \emph{importante} e \emph{água}) com a \textbf{mesma
alocação} dos alvos: duas com três temas e cotas 17/17/16, quatro com dois temas e cotas
25/25. Se os controles tivessem outro número de temas ou de ocorrências, uma diferença entre
alvos e controles na P4 poderia vir só do desenho.

\paragraph{4. Multitema (P4).} 150 palavras de conteúdo com pelo menos 5 ocorrências em pelo
menos 3 temas fora do núcleo, completadas até 20 ocorrências no total (contando as que já
estão na amostra, no núcleo ou nos estratos anteriores), sorteadas em rodízio entre os temas.
Por isso o estrato multitema (\nVerticesMultitema{} vértices) fica abaixo de $150\times20$.
Elas ampliam a P4 para além das 18 palavras escolhidas à mão.

\begin{table}[htbp]
\centering
\caption{Cotas por tema de sentido do desenho final (\cod{configs/experiment.yaml}). São
parâmetros do desenho; as quantidades efetivamente obtidas estão na
Tabela~\ref{tab:amostra-cotas}.}\label{tab:cotas}
\small
\begin{tabular}{@{}l l@{\hspace{2.5em}}l l@{}}
\toprule
\textbf{Alvo} & \textbf{Cotas} & \textbf{Controle} & \textbf{Cotas} \\
\midrule
banco & Economia 17, Computação 17, Geografia 16 & ano & Economia 17, Esporte 17, Música 16 \\
campo & Física 17, Esporte 17, Geografia 16 & século & Física 17, Geografia 17, Política 16 \\
órgão & Biologia 17, Música 17, Política 16 & grande & Física 25, Biologia 25 \\
estado & Física 25, Política 25 & primeiro & Esporte 25, Computação 25 \\
carga & Física 25, Economia 25 & importante & Biologia 25, Economia 25 \\
rede & Computação 25, Geografia 25 & água & Biologia 25, Geografia 25 \\
nota & Música 25, Economia 25 & & \\
capital & Economia 25, Política 25 & & \\
família & Biologia 25, Política 25 & & \\
massa & Física 25, Geografia 25 & & \\
matriz & Biologia 25, Computação 25 & & \\
planta & Biologia 25, Economia 25 & & \\
\bottomrule
\end{tabular}
\end{table}

\paragraph{Sequências.} Cada parágrafo tocado pela amostra vira uma sequência
\tok{<|endoftext|>} $+$ tokens do parágrafo, truncada logo depois da última posição escolhida
nele (Seção~\ref{sec:mod-sequencias}). Nos estratos esparsos, os tokens antes da posição
sorteada são só contexto. A tabela \cod{occurrences.csv} tem uma linha por vértice, ordenada
por (sequência, posição), com origem (tema, artigo, revisão, parágrafo, sentença), posição e
faixa de posição, grupo de prefixo, token (id, texto, \emph{offsets}, vizinhos na sequência),
palavra e categoria, frequências e faixas na amostra e no corpus, estrato, palavra-alvo, tema
de sentido e contexto local. A ordem das linhas define o desempate por posição
(Seção~\ref{sec:red-empates}). Para a separação por sentido na P4, cerca de 20 ocorrências por
alvo são anotadas à mão em \cod{senses.csv}.

\paragraph{A amostra obtida.} Na última execução, a amostra tem \nVertices{} vértices em
\nSequencias{} sequências (\nTokensProcessados{} tokens a processar no modelo) e \nTipos{}
tipos distintos: \nVerticesNucleo{} no núcleo (\nParagrafosNucleo{} parágrafos),
\nVerticesAlvo{} de alvos, \nVerticesControle{} de controles e \nVerticesMultitema{}
multitema. Ficaram \nAlvosMantidos{} alvos na análise.

\tabelaopcional[etapa \etapa{report} a partir de \etapa{sample}]{amostra_composicao}
    {Composição da amostra por estrato e tema (vértices, parágrafos e tipos).}
    {tab:amostra-composicao}
\tabelaopcional[etapa \etapa{report} a partir de \etapa{sample}]{amostra_faixas}
    {Vértices por faixa de frequência (na amostra e no corpus) e por categoria de token.}
    {tab:amostra-faixas}
\tabelaopcional[etapa \etapa{report} a partir de \etapa{sample}]{amostra_cotas}
    {Ocorrências por palavra-alvo e controle em cada tema de sentido: cota, ocorrências que já
    estavam no núcleo, sorteadas fora dele e total obtido; alvos descartados por falta de
    ocorrências.}
    {tab:amostra-cotas}
\figuraopcional[etapa \etapa{report} a partir de \etapa{sample}]{amostra_composicao}
    {Composição da amostra: vértices por estrato, tema, faixa de frequência e categoria de
    token.}
    {fig:amostra-composicao}

% ============================================================
\section{Modelo e extração das representações}\label{sec:modelo}
% ============================================================

% ------------------------------------------------------------
\subsection{Escolha do modelo}\label{sec:mod-escolha}

O modelo analisado é o \textbf{Qwen3-4B-Base}~\cite{qwen3}, na revisão fixada
\cod{\seqsplit{906bfd4b4dc7f14ee4320094d8b41684abff8539}} do Hugging Face. É um decodificador causal com
36 blocos, dimensão oculta $d=2560$, atenção com 32 cabeças de consulta e 8 de chave e valor,
RMSNorm~\cite{zhang2019} e RoPE, cerca de 4 bilhões de parâmetros e pesos em
\emph{bfloat16} (cerca de 7,5 GiB). O tokenizer tem 151.669 ids; a matriz de
\emph{embeddings} tem 151.936 linhas, e as 267 linhas extras são enchimento que o tokenizer
nunca produz.

Quatro candidatos foram considerados:
\begin{description}
    \item[Qwen3-4B pós-treinado.] É a versão para conversa (ajuste fino supervisionado e
    aprendizado por reforço, com modos de raciocínio). O pós-treino adapta as representações ao
    formato de diálogo e a instruções, o que não tem relação com prosa enciclopédica. A
    pergunta do projeto é sobre como um modelo de linguagem reorganiza tokens em contexto, e o
    modelo \emph{Base} é o modelo de linguagem puro, treinado só para prever o próximo token.
    \item[Tucano-2b4.] Treinado em português, com o melhor tokenizer para o nosso texto (1,33
    tokens por palavra e 54\% dos vértices como palavras inteiras, contra 1,78 e 28\% do Qwen3;
    Seção~\ref{sec:viabilidade}). Mas é menor, menos usado e menos representativo dos modelos
    atuais.
    \item[GPT-2 small português.] O modelo do piloto: GPT-2 small (12 blocos, $d=768$) ajustado
    na Wikipédia em português. Tem o tokenizer mais amigável (72\% de palavras inteiras), mas
    é pequeno e antigo, e foi treinado justamente no texto que vamos analisar.
    \item[Qwen3-4B-Base.] Arquitetura atual, multilíngue, com profundidade suficiente para uma
    curva camada a camada (36 blocos) e tamanho dentro da faixa de 4 a 8 bilhões de parâmetros
    prevista na proposta. Cabe na GPU disponível (RTX 4070 de 8 GiB) com parte dos blocos na
    CPU.
\end{description}

A escolha foi o Qwen3-4B-Base. O custo é a fragmentação do português
(Seção~\ref{sec:prob-fragmentacao}), tratado de três formas: as medidas são sempre reportadas
por categoria de token, as palavras-alvo foram restritas às que são um token só, e as
interpretações semânticas se concentram em palavras inteiras de conteúdo.

\paragraph{Embeddings amarrados.} O \cod{config.json} da revisão fixada tem
\cod{tie\_word\_embeddings: true}: a matriz de entrada $E$ também é a matriz de saída, e os
\emph{logits} do próximo token são $E\,\mathrm{norm}(h^{(36)})$~\cite{press2017}. Isso
significa que $e_t$ é, ao mesmo tempo, o vetor com que o token $t$ entra no modelo e a direção
que o modelo usa para prevê-lo. A interpretação da rede lexical e da extensão ``vizinhos no
vocabulário'' (Seção~\ref{sec:res-lente}) leva isso em conta: só a saída normalizada do bloco
36 vive exatamente no espaço que a matriz de saída lê.

% ------------------------------------------------------------
\subsection{Camadas e representações}

Para cada ocorrência são guardados:
\begin{itemize}
    \item \textbf{Lexical} (\cod{lex}): $e_{t_i}$, a linha da matriz de \emph{embeddings} do
    tipo. Os vetores lexicais por ocorrência nunca são materializados: a matriz inteira é
    guardada uma vez (151.669 linhas em bf16, cerca de 780 MB) e $e_{t_i}$ é lido pelo índice.
    \item \textbf{Blocos 1, 18 e 36} (\cod{L01}, \cod{L18}, \cod{L36}): a saída \emph{bruta}
    do bloco, isto é, o fluxo residual depois do bloco $\ell$, sem normalização. O bloco 1 é a
    primeira mistura com o contexto, o 18 ($=\lfloor 36/2\rfloor$) é o intermediário e o 36 é
    o último.
    \item \textbf{Bloco 36 normalizado} (\cod{L36n}): a saída da RMSNorm final, a mesma que os
    \emph{logits} leem. Entra como variante de robustez.
    \item \textbf{Todos os 36 blocos}, gravados num \emph{memmap} em disco (cerca de 2,8 GB
    em bf16 para 15 mil vértices), para a curva camada a camada da P2.
    \item \textbf{Previsão do próximo token}: o argmax de $E\,\mathrm{norm}(h^{(36)}_i)$, usado
    como rótulo na P3 e na checagem da extensão.
\end{itemize}

\paragraph{Por que a saída bruta, e não a normalizada.} Os blocos 1 e 18 só existem na forma
bruta (a RMSNorm final só é aplicada depois do último bloco). Para que a transição
$18\to36$ compare objetos do mesmo tipo, o bloco 36 também é tomado na forma bruta. A RMSNorm
divide cada vetor pela sua raiz da média dos quadrados e multiplica cada dimensão por um ganho
aprendido; o cosseno é invariante à primeira operação, mas não à segunda, então as duas
versões do bloco 36 podem ter vizinhanças diferentes, e essa diferença é medida na robustez.

\paragraph{A pegadinha do \texttt{hidden\_states[36]}.} Por convenção, a tupla
\cod{hidden\_states} devolvida com \cod{output\_hidden\_states=True} tem 37 elementos: o
\emph{embedding} de entrada e a saída de cada bloco. No transformers 5.16, porém, o último
elemento é substituído pela saída \emph{depois} da RMSNorm final: \cod{hidden\_states[36]} não
é a saída bruta do bloco 36 (Seção~\ref{sec:prob-hidden}). Por isso as saídas brutas são
capturadas por \emph{hooks} de \emph{forward} em \cod{layers[0]}, \cod{layers[17]} e
\cod{layers[35]}, e a versão normalizada por um \emph{hook} em \cod{norm}. Hooks em todos os
36 blocos alimentam o \emph{memmap}.

\paragraph{Carregamento.} O modelo é carregado com \cod{AutoModel} (o decodificador sem a cabeça
de linguagem, desnecessária porque a matriz de saída é a própria $E$), com
\cod{dtype=torch.bfloat16} (o argumento \cod{torch\_dtype} está obsoleto no transformers 5),
\cod{device\_map="auto"} e atenção \cod{sdpa}. A execução principal usou uma única RTX A5500
(24 GiB), com \cod{max\_memory} de 20 GiB na GPU: os pesos (7,5 GiB) cabem inteiros e nenhum
bloco fica na CPU (pico de 7,6 GiB alocados, segundo o manifesto da etapa). Numa GPU de 8 GiB,
como a RTX 4070 usada no desenvolvimento, o limite de 6 GiB faz o \cod{accelerate} deixar
parte dos blocos na CPU e transferi-los para a GPU durante o \emph{forward}; isso deixa a
extração mais lenta, mas não muda o resultado, porque cada bloco faz as mesmas contas em bf16.

% ------------------------------------------------------------
\subsection{Sequências, prefixo e grupos de prefixo}\label{sec:mod-sequencias}

\paragraph{Um parágrafo por \emph{forward pass}.} Cada parágrafo é uma sequência processada
sozinha (lote 1), sem \emph{padding}. Assim, o contexto de cada ocorrência é exatamente o texto
anterior a ela no parágrafo, e não depende de quais outros parágrafos estão no lote.

\paragraph{Prefixo.} Toda sequência começa com \tok{<|endoftext|>} (id 151643), que nunca vira
vértice. O tokenizer do Qwen3 não acrescenta token de início (\cod{add\_bos\_token: false}),
e esse id é o separador de documentos do pré-treino (também é o \cod{bos\_token\_id} do
\cod{config.json}), ou seja, o contexto natural de ``início de documento''. Sem prefixo, o
primeiro token do parágrafo ocuparia a posição 0, que em decodificadores funciona como
``sorvedouro'' de atenção e tem ativações de norma enorme~\cite{sun2024}; ele viraria um
vértice atípico por causa da posição, e não do conteúdo.

\paragraph{Truncamento.} Cada sequência é processada só até o último vértice escolhido nela. Como
o modelo é causal, os tokens posteriores não mudam os estados das posições anteriores, e o
truncamento economiza processamento, principalmente nos estratos esparsos.

\paragraph{Grupos de prefixo.} Vértices cujas sequências têm o \emph{mesmo} prefixo de tokens até
eles (inclusive) têm estados matematicamente iguais em todas as camadas; o caso típico é
``\tok{O}'' ou ``\tok{A}'' no começo de vários parágrafos. Na prática, sequências de
comprimentos diferentes podem dar diferenças numéricas mínimas (os \emph{kernels} de atenção
agrupam as contas de forma diferente). Os vértices são agrupados pelo prefixo, os vetores do
primeiro membro são copiados para os outros e a maior diferença encontrada é registrada como
diagnóstico de ruído numérico (\difMaxPrefixo{} na última execução, em \nGruposPrefixo{}
grupos). Com a cópia, esses vértices ficam exatamente empatados e recebem o mesmo tratamento
de empates que a rede lexical, em vez de uma ordem decidida por ruído.

% ------------------------------------------------------------
\subsection{Verificações antes da extração}\label{sec:mod-verificacao}

Antes da extração completa, \cod{gender-networks extract --verify} roda o modelo real em 8
sequências e checa as premissas acima:
\begin{enumerate}
    \item \cod{len(hidden\_states) == 37};
    \item \cod{hidden\_states[0]} é igual a \cod{embed\_tokens(ids)}: a entrada do primeiro bloco
    é o \emph{embedding} puro (o RoPE age dentro da atenção e não soma nada ao fluxo residual);
    \item o \emph{hook} de \cod{layers[k-1]} é igual a \cod{hidden\_states[k]} para
    $k\in\{1,18,35\}$;
    \item \cod{norm(hook\_36)} é igual a \cod{hidden\_states[36]}, o que confirma que este é o
    estado normalizado;
    \item duas execuções dão resultados idênticos.
\end{enumerate}
As quatro primeiras bloqueiam a extração se falharem; a última só gera um aviso.

\tabelaopcional[etapa \etapa{report} a partir de \etapa{extract} (\cod{verify.json})]
    {extracao_verificacao}
    {Resultado das verificações no modelo real: maior diferença absoluta e tolerância de cada
    checagem.}
    {tab:extracao-verificacao}

% ------------------------------------------------------------
\subsection{Diagnósticos das representações}\label{sec:mod-diagnosticos}

Para cada representação, a etapa \etapa{extract} grava em \cod{diagnostics.json}:
\begin{itemize}
    \item quantis da norma dos vetores e normas por faixa de posição (as primeiras posições
    costumam ter normas atípicas);
    \item o cosseno médio entre pares aleatórios de vértices, uma medida de
    \textbf{anisotropia}: se todos os vetores apontam para uma direção comum, o cosseno médio
    é alto e o k-NN por cosseno fica dominado por essa direção~\cite{ethayarajh2019};
    \item as dimensões dominantes (\emph{ativações massivas}): poucas coordenadas com valores
    muito maiores que as demais, conhecidas em modelos grandes~\cite{sun2024,timkey2021}, que
    podem dominar o cosseno;
    \item o perfil camada a camada (norma média e cosseno médio de pares aleatórios nos 36
    blocos), que antecipa onde a geometria do espaço muda.
\end{itemize}
A previsão do próximo token (\cod{pred\_next.npy}) é calculada como os \emph{logits} do modelo,
$E\,\mathrm{norm}(h^{(36)}_i)$, sem carregar a cabeça de linguagem, graças aos
\emph{embeddings} amarrados.
Esses diagnósticos motivam duas variantes de robustez: o cosseno \emph{centrado} (vetores menos
a média sobre as ocorrências, o que remove a direção comum) e o k-NN mútuo, menos sensível a
\emph{hubs} (Seção~\ref{sec:red-robustez}).

\tabelaopcional[etapa \etapa{report} a partir de \etapa{extract}]{extracao_diagnosticos}
    {Diagnósticos por representação: quantis de norma, cosseno médio de pares aleatórios e
    dimensões dominantes.}
    {tab:extracao-diagnosticos}
\figuraopcional[etapa \etapa{report} a partir de \etapa{extract}]{extracao_camadas}
    {Norma média e cosseno médio entre pares aleatórios ao longo dos 36 blocos.}
    {fig:extracao-camadas}
\figuraopcional[etapa \etapa{report} a partir de \etapa{extract}]{extracao_dimensoes}
    {Dimensões dominantes: magnitude das maiores coordenadas por camada e norma por faixa de
    posição.}
    {fig:extracao-dimensoes}

% ============================================================
\section{Construção das redes}\label{sec:redes}
% ============================================================

Todas as redes do experimento são grafos de $k$ vizinhos mais próximos (k-NN) sob o cosseno.
Esta seção descreve os conjuntos de vértices, o cálculo das similaridades, os três tratamentos
de empate e as versões direcionada, por união e mútua. A construção é feita pela etapa
\etapa{knn}.

% ------------------------------------------------------------
\subsection{Conjuntos de vértices}\label{sec:red-vertices}

Os tipos da amostra (os \cod{token\_id} distintos entre as ocorrências, cerca de 3,3 mil) saem
só da amostragem, antes de qualquer k-NN. O k-NN roda então sobre três conjuntos de vértices:
\begin{description}
    \item[As $\approx$15 mil ocorrências da amostra.] As redes principais, onde P1--P4 são
    respondidas: lexical por ocorrência (\cod{lex}) e contextuais (\cod{L01}, \cod{L18},
    \cod{L36}), mais as variantes de robustez (\cod{L36n} e as centradas \cod{L01c},
    \cod{L18c}, \cod{L36c}). Elas se restringem à amostra porque só tokens que aparecem num
    texto têm vetor contextual. Todas têm exatamente os mesmos \nVertices{} vértices.
    \item[O vocabulário inteiro.] A rede lexical de tipos (\cod{vocab}), com os 151.669 ids do
    tokenizer menos os 26 especiais, ou seja, 151.643 vértices (\nTiposVocab{} na última
    execução). As 267 linhas de enchimento da matriz de \emph{embeddings} ficam de fora porque
    nenhum texto as produz. Cada vértice tem como atributos a frequência no corpus e na amostra
    (zero para a maioria), a classe de escrita e se aparece no corpus (\nTiposVocabCorpus{}
    tipos).
    \item[Os $\approx$3,3 mil tipos da amostra.] A tabela auxiliar de vizinhos, não reportada
    como rede, usada na derivação exata da rede lexical por ocorrência e no $T_i$ lexical da
    variante (c) (Seções~\ref{sec:vocab-vs-amostra} e~\ref{sec:red-derivacao}).
\end{description}

A classe de escrita de cada tipo do vocabulário é uma de: latina, chinês/japonês/coreano (CJK),
mista (código, misturas de escritas), pontuação, outras escritas (cirílico, árabe\ldots), pedaço
de bytes (decodifica para U+FFFD), espaço, dígitos e especial. Uma medição na fase de
planejamento deu 50\% de tokens de escrita latina, 20\% CJK e 13\% mistos: a maior parte do
vocabulário nunca aparece num texto em português. Tokens pouco treinados podem ter
\emph{embeddings} quase iguais e formar aglomerados ou \emph{hubs} artificiais; isso é
diagnosticado e reportado nos resultados da rede do vocabulário.

% ------------------------------------------------------------
\subsection{Similaridade em \emph{float64}}\label{sec:red-cosseno}

A similaridade entre dois vértices é o cosseno
\begin{equation}
    s_{ij} = \frac{x_i^\top x_j}{\lVert x_i\rVert\,\lVert x_j\rVert}.
\end{equation}
Os vetores são guardados em bf16, convertidos exatamente para \emph{float64}, normalizados
linha a linha e multiplicados em blocos de até 1.024 linhas na CPU (limitados a 512~MiB por
bloco; cerca de 440 linhas na rede do vocabulário). O \emph{float64} é necessário
por causa da tolerância de empate $\eps=10^{-6}$: em \emph{float32}, um produto interno com
$d=2560$ termos acumula erro de arredondamento da ordem de $10^{-6}$ a $10^{-5}$ no cosseno,
comparável à própria tolerância, o que criaria empates falsos e desfaria empates verdadeiros.
Em \emph{float64} o erro fica abaixo de $10^{-12}$.

A mesma rotina serve para a rede do vocabulário: 151,6 mil vetores normalizados ocupam cerca de
3,1 GB em \emph{float64}, e o cálculo em blocos mantém a memória sob controle.

% ------------------------------------------------------------
\subsection{Listas de candidatos}\label{sec:red-candidatos}

Para cada representação, a etapa guarda para cada vértice $i$ uma lista de candidatos ordenada
por similaridade decrescente (e índice crescente), sem o próprio $i$, em formato CSR
(\cod{cand\_<rep>.npz}). A lista tem os $K=128$ mais similares, com uma garantia: ela contém
\emph{todo} $j$ com $s_{ij}\ge s_{i(k_{\max})}-\eps$, onde $s_{i(k_{\max})}$ é a
$k_{\max}$-ésima maior similaridade da linha e $k_{\max}=20$ é o maior $k$ usado. Se um bloco de
empates passar do $K$-ésimo candidato, a linha inteira é recalculada e todos os empatados
entram. Com essa garantia, todos os valores de $k\le k_{\max}$ e todos os tratamentos de empate
saem da mesma lista, sem recalcular similaridades.

% ------------------------------------------------------------
\subsection{Rede direcionada, por união e mútua}\label{sec:red-versoes}

A relação de $k$ vizinhos mais próximos não é simétrica: $j\in\N_i$ não implica $i\in\N_j$. O
objeto produzido pelo k-NN é, portanto, um \textbf{grafo direcionado}
\begin{equation}
    G^{\rightarrow} = \bigl(V,\ \{\,i\to j : j\in\N_i\,\}\bigr),
\end{equation}
em que todo vértice tem grau de saída $k$ (nas variantes de tamanho fixo) e o grau de entrada
varia: vértices escolhidos por muitos outros são \emph{hubs}. A versão direcionada é usada nas
medidas de vizinhança (Jaccard e dominância, que dependem de $\N_i$), na distribuição de graus
de entrada e na reciprocidade.

Para clusterização, distâncias, diâmetro, componentes e comunidades, usamos a versão
\textbf{não direcionada por união},
\begin{equation}
    E^{\cup} = \bigl\{\{i,j\} : j\in\N_i \ \text{ou}\ i\in\N_j\bigr\},
\end{equation}
em que todo vértice tem grau pelo menos $k$ e, portanto, não há vértices isolados (o que não
garante conexidade, como se vê abaixo). A versão
\textbf{mútua},
\begin{equation}
    E^{\cap} = \bigl\{\{i,j\} : j\in\N_i \ \text{e}\ i\in\N_j\bigr\},
\end{equation}
é mais conservadora (só relações confirmadas pelos dois lados, o que corta as arestas que
chegam a \emph{hubs}) e entra como robustez. Distância média e diâmetro são calculados no maior
componente conexo, porque as duas versões podem se fragmentar. A mútua, por construção; a
união, por um mecanismo específico destas redes: um conjunto de mais de $k$ vértices com vetores
idênticos (as ocorrências de um tipo com $f_t>k$ na rede lexical, ou um grupo de prefixo com
mais de $k$ ocorrências nas camadas contextuais) manda todas as suas arestas de saída para
dentro de si e forma um componente fechado se nenhum vértice de fora escolher algum deles. Na
rede lexical isso atinge justamente os tipos frequentes do desenho, e as distâncias do maior
componente os excluem; por isso a tabela de métricas informa, para cada representação, o
número de componentes e a fração de vértices no maior (Seção~\ref{sec:met-estruturais}).

% ------------------------------------------------------------
\subsection{Tratamento de empates}\label{sec:red-empates}

\paragraph{Por que os empates importam.} Na rede lexical por ocorrência, as $f_{t_i}-1$ outras
ocorrências do tipo de $i$ têm $s_{ij}=1$ exatamente. Se $f_{t_i}-1>k$, a escolha de quais delas
entram em $\N_i$ é arbitrária. Se $f_{t_i}-1<k$, as vagas restantes vão para ocorrências do
tipo mais próximo, que também são idênticas entre si e, portanto, também empatam. Empates exatos
também aparecem nas camadas contextuais, nos grupos de prefixo (Seção~\ref{sec:mod-sequencias}).
Como as redes contextuais são comparadas com a lexical, o tratamento de empates afeta
diretamente as medidas de reorganização.

\paragraph{Definição comum.} Seja $s_{i(k)}$ a $k$-ésima maior similaridade da linha $i$. Os
candidatos se dividem em
\begin{equation}
    F_i = \{\,j : s_{ij} > s_{i(k)}+\eps\,\},\qquad
    B_i = \{\,j : |s_{ij}-s_{i(k)}|\le\eps\,\},\qquad
    r_i = k-|F_i|,
\end{equation}
isto é, os que estão com certeza entre os $k$ primeiros ($F_i$, de ``\emph{front}'') e o bloco
empatado na fronteira ($B_i$, de ``\emph{boundary}''). Sempre vale $|F_i|<k\le|F_i|+|B_i|$, e
faltam $r_i$ vizinhos a escolher em $B_i$. Usamos $\eps=10^{-6}$ na configuração principal e
$\eps=10^{-4}$ como teste de sensibilidade. As três variantes são:

\begin{enumerate}[label=(\alph*)]
    \item \textbf{Sorteio com sementes (principal).} $\N_i=F_i\cup R_i$, com $R_i\subset B_i$
    sorteado uniformemente sem reposição, $|R_i|=r_i$. A etapa \etapa{knn} grava o sorteio com
    $R=20$ sementes para toda representação, e as medidas de vizinhança ($J$, $D$, piso de
    ruído) usam todas elas. As métricas estruturais usam as 20 sementes só na rede lexical, onde
    os empates são maciços, e são reportadas como média $\pm$ desvio entre elas; nas camadas
    contextuais os empates vêm só de grupos de prefixo (vetores idênticos), o sorteio quase não
    muda a rede, e as métricas estruturais e as comunidades usam a semente 0, com o desempate
    por posição como checagem. As comunidades da rede lexical usam as sementes 0 e 1 (a segunda
    mostra quanto a partição depende do sorteio). A semente $r$ da
    representação \cod{rep} com $k$ vizinhos usa o gerador
    \cod{default\_rng([438, crc32(rep), k, r])}, de modo que cada combinação tem um fluxo
    aleatório próprio e reprodutível, que não depende da ordem em que as redes são calculadas.
    O desempate pela \textbf{posição} (os $r_i$ menores \cod{occurrence\_id} de $B_i$) é mantido
    como referência determinística; com ele e a versão por união, o k-NN novo reproduz
    exatamente o \cod{build\_union\_knn\_graph} do piloto (teste de regressão).
    \item \textbf{Todos os empatados.} $\N_i=F_i\cup B_i$. Elimina a arbitrariedade, ao custo de
    um grau de saída variável ($|\N_i|\ge k$). Robustez.
    \item \textbf{Tipos distintos.} As ocorrências do próprio tipo são ignoradas, e a vizinhança
    passa a ser um conjunto $T_i$ de $k$ \emph{tipos} $u\neq t_i$. Cada tipo recebe a nota da
    sua ocorrência mais similar a $i$,
    \begin{equation}
        \sigma_{iu} = \max_{j:\,t_j=u} s_{ij},
    \end{equation}
    e $T_i$ são os $k$ tipos de maior nota, com a mesma regra de $F$ e $B$ sobre $\sigma$
    (desempate pelo menor \cod{token\_id} ou por sorteio). Nas camadas contextuais, $T_i$ diz
    \emph{de quais outros tokens} a ocorrência se aproxima no seu contexto. A variante entra só
    nas medidas de vizinhança (Jaccard sobre conjuntos de tipos), porque a dominância lexical
    não faz sentido quando o próprio tipo é excluído.
\end{enumerate}

\paragraph{Dois exemplos na rede lexical.} Com $k=10$:
\begin{itemize}
    \item um tipo com $f_t=50$: $F_i=\emptyset$ e $B_i$ são as 49 outras ocorrências do tipo
    ($s=1$); $r_i=10$. Com o desempate por posição, as 50 ocorrências escolhem as mesmas 10 de
    menor índice, que viram \emph{hubs} com grau de entrada perto de 50, um artefato puro. Com
    o sorteio, cada ocorrência escolhe 10 das 49 ao acaso, e os graus de entrada se espalham em
    torno de 10. Foi esse artefato, visível já no piloto com 104 tokens, que motivou o sorteio
    como regra principal;
    \item um tipo com $f_t=4$ cujo tipo mais próximo $u$ tem $f_u\ge7$: $F_i$ são as 3 outras
    ocorrências do tipo ($s=1$), $B_i$ são as $f_u$ ocorrências de $u$ (todas com a mesma
    similaridade) e $r_i=7$.
\end{itemize}

\tabelaopcional[etapa \etapa{report} a partir de \etapa{knn} (\cod{\_manifest.json})]{knn_empates}
    {Estatísticas de empates por representação e $k$: linhas com empate na fronteira, maior
    bloco empatado, linhas recalculadas por inteiro e tempo de cálculo.}
    {tab:knn-empates}

% ------------------------------------------------------------
\subsection{Derivação exata da rede lexical por ocorrência}\label{sec:red-derivacao}

Na rede lexical, $x_i=e_{t_i}$, então $s_{ij}=c(t_i,t_j)$, onde $c(t,u)$ é o cosseno entre os
\emph{embeddings} dos tipos $t$ e $u$: a similaridade entre duas ocorrências só depende dos
seus tipos. A linha ordenada de $i$ é, portanto:
\begin{enumerate}
    \item as $f_{t_i}-1$ outras ocorrências do mesmo tipo, com $s=1$;
    \item para cada tipo $u$ da amostra, em ordem decrescente de $c(t_i,u)$, as $f_u$
    ocorrências de $u$, todas com a mesma similaridade.
\end{enumerate}
Os candidatos de $i$ (e $\N_i$ em qualquer variante) saem exatamente da matriz completa de
similaridade entre os cerca de 3,3 mil tipos da amostra ($U\times U$, barata nesse tamanho) e
das contagens $f_u$, sem calcular a matriz $15\,000\times15\,000$. Como há um único valor de
similaridade por par de tipos, os blocos de empate são exatos, sem depender de arredondamento.
Um teste compara a derivação com a força bruta. (O parâmetro \texttt{type\_candidates}${}=64$
é outra coisa: o número de tipos candidatos guardados por linha para a variante (c).)

Na variante (c), a nota de um tipo é $\sigma_{iu}=c(t_i,u)$, e $T_i$ é simplesmente o conjunto
dos $k$ tipos da amostra mais próximos de $t_i$ na tabela auxiliar. Quase nunca há empate na
fronteira: empates exatos só entre tipos com \emph{embeddings} idênticos e, com a tolerância
$\eps$, raros casos de tipos distintos com cossenos a menos de $\eps$ um do outro (numa
aproximação com os 3,3 mil tipos mais frequentes da viabilidade, 2 linhas com $k=10$ e 3 com
$k=20$).

% ------------------------------------------------------------
\subsection{Rede lexical de tipos (vocabulário)}\label{sec:red-vocab}

A rede do vocabulário usa os mesmos passos sobre os 151.643 vetores $e_t$: cosseno em
\emph{float64}, candidatos com a garantia de completude e $k\in\{5,10,20\}$, com $k=10$ como
principal. Empates exatos só acontecem entre tipos com \emph{embeddings} idênticos (tokens
nunca treinados, por exemplo): na revisão fixada, 1.903 linhas do vocabulário repetem o vetor
de pelo menos outra linha, em 58 grupos (o maior com 524 tipos). Com a tolerância $\eps$, raros
empates entre tipos distintos também entram em $B$ (cerca de 0,1\% das linhas numa amostra com
$k=10$). Os empates são resolvidos por posição (\cod{token\_id}) e por um sorteio. Um grupo de
vetores idênticos com mais de $k$ tipos manda todas as suas arestas de saída para dentro de si
e pode formar um componente fechado, pelo mesmo mecanismo da
Seção~\ref{sec:red-versoes}. As métricas estruturais usam a versão por união; a
distribuição de graus de entrada, a direcionada.

% ------------------------------------------------------------
\subsection{Configuração principal e grade de robustez}\label{sec:red-robustez}

A configuração principal é fixa: $k=10$, versão por união (direcionada para vizinhanças e graus
de entrada), empates pela variante (a) (20 sementes na rede lexical e nas medidas de
vizinhança, semente 0 nas métricas estruturais das camadas contextuais), $\eps=10^{-6}$ e saídas
brutas dos
blocos. A Tabela~\ref{tab:grade} lista as variações da robustez; os resultados estão na
Seção~\ref{sec:robustez}.

\begin{table}[htbp]
\centering
\caption{Configuração principal e variações de robustez.}\label{tab:grade}
\small
\begin{tabularx}{\linewidth}{@{}l l X@{}}
\toprule
\textbf{Dimensão} & \textbf{Principal} & \textbf{Robustez e motivo} \\
\midrule
$k$ & 10 & 5 e 20: a estrutura não deve depender de um $k$ específico \\
Simetrização & união & mútua: remove arestas que só existem por causa de \emph{hubs} \\
Empates & (a) sorteio, 20 sementes & posição (referência do piloto); (b) todos os empatados;
(c) tipos distintos \\
Tolerância & $\eps=10^{-6}$ & $\eps=10^{-4}$: sensibilidade da definição de empate \\
Bloco 36 & saída bruta (\cod{L36}) & normalizada (\cod{L36n}), o espaço que os
\emph{logits} leem \\
Centragem & não & \cod{L01c}, \cod{L18c}, \cod{L36c}: vetores menos a média sobre as
ocorrências, que remove a direção comum (anisotropia, ativações massivas) \\
Comunidades & Leiden, resolução 1 & resoluções 0,5 e 2; Louvain como checagem \\
Vértices & amostra inteira & só o núcleo ($\approx$11 mil vértices, ainda acima de $10^4$);
planejado: exige um $k$-NN próprio, ainda sem etapa \\
\bottomrule
\end{tabularx}
\end{table}

% ============================================================
\section{Métricas e métodos de análise}\label{sec:metodos}
% ============================================================

% ------------------------------------------------------------
\subsection{Métricas estruturais}\label{sec:met-estruturais}

Em todas as redes da configuração principal calculamos as métricas exigidas pela disciplina,
mais reciprocidade e \emph{hubs}. Seja $G=(V,E)$ a rede não direcionada por união, $n=|V|$,
$m=|E|$ e $k_i$ o grau do vértice $i$.
\begin{description}
    \item[Tamanho e grau.] $n$, $m$, grau médio $\langle k\rangle=2m/n$ e densidade
    $\rho=2m/[n(n-1)]$. Nas redes k-NN, $m$ fica entre $nk/2$ (todas as escolhas recíprocas) e
    $nk$ (nenhuma recíproca), então $\langle k\rangle\in[k,2k]$ e a densidade é muito baixa
    por construção; o interesse está em como $m$ varia entre camadas, que reflete a
    reciprocidade.
    \item[Distribuição de graus.] Na rede não direcionada, e de entrada e de saída na
    direcionada. Na variante (a), o grau de saída é $k$ para todos, e toda a informação está no
    grau de entrada $k^{\mathrm{in}}_i=|\{j: i\in\N_j\}|$, cuja média também é $k$. Reportamos a
    CCDF $P(K^{\mathrm{in}}\ge x)$ e a assimetria (\emph{skewness}) de $k^{\mathrm{in}}$, a medida
    usual de \emph{hubness}~\cite{radovanovic2010}: assimetria alta significa poucos vértices
    escolhidos por muitos.
    \item[Reciprocidade.] Fração das arestas direcionadas $i\to j$ para as quais $j\to i$
    também existe.
    \item[Clusterização.] Global (transitividade),
    $C=3\times(\text{triângulos})/(\text{trios conexos})$, e local média,
    \begin{equation}
        \bar C=\frac{1}{|V_{\ge2}|}\sum_{i\in V_{\ge2}} C_i,\qquad C_i=\frac{2\,e_i}{k_i(k_i-1)},
    \end{equation}
    onde $e_i$ é o número de arestas entre os vizinhos de $i$. $C_i$ não é definido quando
    $k_i<2$, e esses vértices ficam fora da média: $V_{\ge2}$ é o conjunto dos vértices de
    grau $\ge 2$ (na rede por união $k_i\ge k$, então $V_{\ge2}=V$).
    \item[Distâncias.] Distância média $\langle\ell\rangle$ entre pares do maior componente
    conexo, diâmetro e histograma das distâncias, por busca em largura exata a partir de todos
    os vértices.
    \item[Componentes.] Número de componentes conexos, fração de vértices no maior componente
    e distribuição dos tamanhos, para toda representação e nas versões por união e mútua. A
    versão por união \emph{não} é necessariamente conexa: um tipo com $f_t>k$ manda todas as
    $k$ arestas de saída de cada ocorrência para outras ocorrências do próprio tipo ($s=1$) e
    vira um componente fechado se nenhum vértice de outro tipo escolher alguma delas; o mesmo
    vale, nas camadas contextuais, para um grupo de prefixo com mais de $k$ ocorrências
    (vetores idênticos) e, na rede do vocabulário, para um grupo de mais de $k$ tipos com
    \emph{embeddings} idênticos (há 58 grupos desses, o maior com 524 tipos;
    Seção~\ref{sec:red-vocab}). Espera-se, portanto, que a rede lexical por união tenha vários
    componentes pequenos formados justamente pelos tipos frequentes do desenho (alvos,
    controles e multitemáticos); numa simulação com o perfil de frequências da amostra híbrida,
    $k=10$ e vetores aleatórios, foram 23 a 24 componentes, com o maior reunindo cerca de 96\%
    dos vértices. Como distância média e diâmetro são calculados no maior componente, na rede
    lexical eles excluem esses tipos, e a comparação de $\langle\ell\rangle$ entre
    representações deve ser lida junto com a fração do maior componente, que a tabela de
    métricas da rede por união informa.
\end{description}

As métricas são calculadas com o igraph~\cite{csardi2006}, cuja busca em largura roda em C:
as distâncias exatas de uma rede de 15 mil vértices levam segundos, contra 20 a 25 minutos por
rede no NetworkX do piloto. O NetworkX fica como referência nos testes (as duas bibliotecas
devem dar os mesmos valores em grafos pequenos). Na rede do vocabulário (151,6 mil vértices),
as distâncias exatas são calculadas só na configuração principal (estimativa: de dezenas de
minutos a cerca de uma hora); nas variantes de robustez, a distância média é estimada por busca
em largura a partir de 500 fontes sorteadas, e o diâmetro, por um limite inferior de varredura
dupla (a partir de um vértice, ir ao mais distante e repetir).

% ------------------------------------------------------------
\subsection{Persistência da vizinhança: Jaccard}\label{sec:met-jaccard}

A reorganização da vizinhança de uma ocorrência entre as redes $a$ e $b$ (por exemplo, lexical
e bloco 1) é medida pela similaridade de Jaccard
\begin{equation}
    J_i(a,b)=\frac{|\N_i^{(a)}\cap\N_i^{(b)}|}{|\N_i^{(a)}\cup\N_i^{(b)}|}.
\end{equation}
Com vizinhanças de tamanho fixo $k$ (variante (a)) e $c=|\N_i^{(a)}\cap\N_i^{(b)}|$ vizinhos em
comum, $J_i=c/(2k-c)$: $J=1$ quando nada muda e $J=0$ quando nenhum vizinho se mantém. Na
variante (b), os conjuntos têm tamanhos diferentes e a fórmula geral vale; na variante (c), a
mesma fórmula é aplicada aos conjuntos de tipos $T_i$. Chamamos de \emph{mudança} o complemento
$M_i=1-J_i$. Com 20 sementes na rede lexical, $J_i(\mathrm{lex},\ell)$ é calculado para cada
semente e resumido por média e desvio.

% ------------------------------------------------------------
\subsection{Dominância lexical e seu teto}\label{sec:met-dominancia}

A dominância lexical é a fração de vizinhos do mesmo tipo,
\begin{equation}
    D_i=\frac{|\{\,j\in\N_i : t_j=t_i\,\}|}{|\N_i|},
\end{equation}
com $|\N_i|=k$ na variante (a) e o tamanho variável na variante (b). Na rede lexical com a
variante (a), as outras ocorrências do tipo são as mais similares ($s=1$), então
\begin{equation}
    D_i^{\mathrm{lex}}=D^{\max}_i=\frac{\min(f_{t_i}-1,\,k)}{k}
\end{equation}
exatamente (um teste verifica isso). Esse teto depende de $f_t$: uma ocorrência de um tipo com
$f_t=4$ nunca passa de $D=0{,}3$ com $k=10$. Para comparar faixas de frequência, reportamos
também a dominância relativa ao teto, $\tilde D_i=D_i/D^{\max}_i$ (definida para $f_t\ge2$ e
analisada para $f_t\ge3$). A queda $\Delta D$ entre duas camadas mede quanto a identidade do
token perde espaço na vizinhança.

% ------------------------------------------------------------
\subsection{Piso de ruído}\label{sec:met-ruido}

Na rede lexical, parte dos vizinhos é sorteada entre candidatos empatados. Duas construções com
sementes diferentes, $r\neq r'$, já dão $J_i(\mathrm{lex}_r,\mathrm{lex}_{r'})<1$ sem nenhuma
mudança de representação. O \textbf{piso de ruído} é essa mudança produzida só pelo desempate,
\begin{equation}
    M^{\text{ruído}}_i = 1-\frac{1}{m}\sum_{q=0}^{m-1}
        J_i\bigl(\mathrm{lex}_{2q},\mathrm{lex}_{2q+1}\bigr),\qquad m=\lfloor R/2\rfloor,
\end{equation}
isto é, a média sobre os pares \emph{disjuntos} de sementes $(0,1),(2,3),\dots$ (com $R=20$
sementes, 10 pares; com $R$ ímpar a última fica de fora). Usar pares disjuntos, em vez de
todos os $\binom{R}{2}$ pares, mantém as estimativas de cada par independentes entre si; o
valor esperado é o mesmo. O piso é
calculado por ocorrência e resumido por faixa de frequência (ele é maior para tipos frequentes,
em que o bloco empatado é grande). Uma mudança lexical $\to$ contextual só é interpretada como
reorganização se $M_i(\mathrm{lex},\ell)$ ficar acima do piso da faixa correspondente.

% ------------------------------------------------------------
\subsection{Jaccard por composição de tipos}\label{sec:met-jw}

Como complemento que não depende da semente, comparamos a \emph{composição por tipos} das
vizinhanças. Seja $n^{(a)}_{iu}$ o número de vizinhos de $i$ do tipo $u$ na rede $a$. O Jaccard
ponderado é
\begin{equation}
    J^{w}_i(a,b)=\frac{\sum_u \min\bigl(n^{(a)}_{iu},n^{(b)}_{iu}\bigr)}
                      {\sum_u \max\bigl(n^{(a)}_{iu},n^{(b)}_{iu}\bigr)}.
\end{equation}
Na rede lexical, o bloco empatado da fronteira costuma ser formado por ocorrências de um único
tipo (as do tipo seguinte no ranking), e então qualquer sorteio dá a mesma composição: $J^w$ é
invariante à semente sempre que o bloco da fronteira tem um só tipo (um teste verifica isso em
dados sintéticos). Isso vale para quase todas as linhas, mas não para todas: dois tipos com
\emph{embeddings} idênticos, ou com cossenos a menos de $\eps$ um do outro, dividem o mesmo
bloco, e nessas linhas a composição varia com a semente. A fração dessas exceções é
reportada junto com $J^w$. Ele responde a uma pergunta um pouco
diferente de $J$: a ocorrência continua perto dos mesmos \emph{tipos}, ainda que de outras
ocorrências deles?

% ------------------------------------------------------------
\subsection{Estratificação}\label{sec:met-estratos}

As distribuições de $J$, $M$, $D$ e $\tilde D$ são reportadas por faixa de frequência (na
amostra e no corpus; Seção~\ref{sec:faixas}), por categoria de token, por estrato da amostra
(núcleo, alvos, controles, multitema) e com controle pela faixa de posição. Na P2, além da
mudança média por transição, contamos a fração de vértices cuja maior mudança acontece em cada
transição (lexical $\to$ 1, $1\to18$, $18\to36$) e traçamos as curvas $J(\ell,\ell+1)$ e
$D(\ell)$ nos 36 blocos, a partir do \emph{memmap} de todas as camadas.

% ------------------------------------------------------------
\subsection{Comunidades}\label{sec:met-comunidades}

As comunidades são detectadas na rede por união com o algoritmo de
Leiden~\cite{traag2019}, que melhora o Louvain~\cite{blondel2008} garantindo comunidades
conexas, maximizando a modularidade~\cite{newman2006}
\begin{equation}
    Q=\frac1{2m}\sum_{ij}\Bigl(A_{ij}-\gamma\frac{k_ik_j}{2m}\Bigr)\,\delta(c_i,c_j),
\end{equation}
com resolução $\gamma=1$ na configuração principal e $\gamma\in\{0{,}5;\,2\}$ na robustez. No
igraph, o objetivo padrão do \cod{community\_leiden} é o CPM (\emph{constant Potts model}),
cuja resolução tem outro significado (com resolução 1, o CPM tende a produzir comunidades
unitárias); por isso o código sempre passa \cod{objective\_function="modularity"}. O Leiden roda
com 10 sementes (o igraph usa o gerador \cod{random} do Python, que é semeado a cada execução);
fica a partição de maior modularidade, e a \textbf{estabilidade} é o NMI médio entre as
partições das 10 execuções. O Louvain serve de checagem.

% ------------------------------------------------------------
\subsection{Concordância entre partições (P3)}\label{sec:met-concordancia}

As comunidades de cada camada são comparadas com partições conhecidas dos vértices: tipo de
token, tema, artigo, parágrafo, sentença, próximo token previsto, próximo token real e faixa de
posição (esta como controle de artefato). Três medidas são usadas:
\begin{description}
    \item[NMI.] Informação mútua normalizada, $\mathrm{NMI}(U,V)=2\,I(U;V)/[H(U)+H(V)]$, via
    \cod{igraph.\allowbreak compare\_communities}. O NMI cresce com o número de classes mesmo entre
    partições sem relação nenhuma~\cite{vinh2010}, o que é grave aqui: ``parágrafo'' tem cerca
    de 80 classes no núcleo e ``tipo'', milhares. Por isso cada NMI é acompanhado de uma
    \textbf{linha de base de permutação}: o NMI médio entre as comunidades e 20 permutações
    aleatórias do rótulo, que preservam o número e o tamanho das classes. Reportamos o NMI
    bruto e $\mathrm{NMI}-\mathrm{NMI}_{\mathrm{perm}}$.
    \item[ARI.] Índice de Rand ajustado~\cite{hubert1985}, já corrigido pelo acaso (zero em
    média para partições independentes).
    \item[Pureza.] $P=\frac1n\sum_c\max_\lambda |c\cap\lambda|$: a fração de vértices que
    pertencem ao rótulo majoritário da sua comunidade. É fácil de interpretar, mas favorece
    comunidades pequenas, então é lida junto com as outras duas.
\end{description}
A P3 é analisada \textbf{só no núcleo}, onde os parágrafos estão inteiros e artigo e parágrafo
são rótulos distintos (dois parágrafos por artigo); nos estratos esparsos, cada parágrafo
contribui com uma ou poucas ocorrências e os rótulos de contexto quase coincidem com o tipo. O
tema também é analisado na rede inteira.

% ------------------------------------------------------------
\subsection{Separação intra-tipo (P4)}\label{sec:met-p4}

Para cada tipo com $f_t\ge10$ (alvos, controles, multitema e palavras funcionais separados):
\begin{itemize}
    \item \textbf{número efetivo de comunidades} ocupadas pelas suas ocorrências,
    $e^{H_t}$ com $H_t=-\sum_c p_{tc}\ln p_{tc}$ e $p_{tc}$ a fração das ocorrências de $t$
    na comunidade $c$ (vale 1 quando todas estão juntas; na rede lexical tende a 1, porque
    ocorrências do mesmo tipo são idênticas e ligadas entre si);
    \item \textbf{NMI entre comunidade e tema} dentro do tipo: as ocorrências se separam
    \emph{segundo o tema}, ou só se espalham?
    \item \textbf{autossimilaridade}: cosseno médio entre as próprias ocorrências do tipo
    (1 na camada lexical)~\cite{ethayarajh2019};
    \item \textbf{diferença entre o cosseno médio dentro do mesmo tema e entre temas};
    \item \textbf{separação pelo sentido anotado} (cerca de 20 ocorrências por alvo, em
    \cod{senses.csv}), que controla o fato de o tema ser só um indicador aproximado do sentido;
    \item \textbf{redes ego} de 2 ou 3 alvos, com layout fixo entre camadas, para inspeção
    qualitativa.
\end{itemize}
A hipótese H4 é testada comparando a distribuição dessas medidas entre alvos e controles,
que têm a mesma alocação de temas e cotas (Seção~\ref{sec:amostra}).

% ------------------------------------------------------------
\subsection{Rede do vocabulário}\label{sec:met-vocab}

Na rede lexical de tipos, além das métricas estruturais: distribuição de graus e \emph{hubs}
(quais tokens são escolhidos por muitos: tokens pouco treinados? pedaços de bytes?),
comunidades comparadas com a classe de escrita (NMI e ARI), a fração de tokens que aparecem no
corpus em cada comunidade e a vizinhança das palavras-alvo no vocabulário inteiro.

% ------------------------------------------------------------
\subsection{Extensão opcional: vizinhos no vocabulário}\label{sec:met-lente}

Para cada ocorrência e representação (lexical, blocos 1, 18 e 36, e bloco 36 normalizado),
calculamos o ranking das linhas da matriz de \emph{embeddings} inteira pelo cosseno (em
\emph{float32} na GPU; aqui não há tratamento de empates, porque só a posição no ranking é
usada) e medimos, camada a camada, a posição do \emph{próprio} token e a do \emph{próximo}
token nesse ranking, e o Jaccard entre os $k=20$ vizinhos no vocabulário de camadas
consecutivas. A pergunta: em que camada a ocorrência deixa de parecer com o próprio token e
passa a parecer com o próximo? Nas camadas intermediárias, a comparação mistura espaços que não
foram treinados para se alinhar; só o bloco 36 normalizado é exatamente o que a matriz de saída
lê. Duas checagens validam a extensão: na camada lexical, o próprio token é o primeiro do
ranking; no bloco 36 normalizado, a linha de \emph{maior produto interno} (os \emph{logits}
recalculados a partir de $E$) coincide com o argmax dos \emph{logits} do próprio modelo em
praticamente todas as ocorrências (\concordanciaPredicao{} na última execução). Pelo cosseno, a
primeira linha pode diferir, porque as linhas de $E$ não têm todas a mesma norma.

% ============================================================
\section{Resultados}\label{sec:resultados}
% ============================================================

Esta seção ainda é só estrutura. Cada subseção diz o que vai conter, quais medidas da
Seção~\ref{sec:metodos} entram e de qual etapa vêm os dados. Os quadros ``pendente'' indicam o
arquivo esperado; quando a etapa \etapa{report} for executada depois das etapas indicadas, eles
são substituídos pelas figuras e tabelas geradas, sem mudar o texto ao redor.

Duas regras de leitura valem para toda a seção (Seção~\ref{sec:perguntas}): nenhuma mudança é
interpretada abaixo do piso de ruído, e nenhum NMI é lido sem a linha de base de permutação.
Salvo indicação, os valores são da configuração principal ($k=10$, união, variante (a)). As
métricas estruturais da rede lexical são a média $\pm$ desvio entre as 20 sementes do
desempate; as das camadas contextuais e todas as comunidades usam a semente 0
(Seção~\ref{sec:red-empates}); as medidas de vizinhança ($J$, $D$) usam as 20 sementes em todas
as representações.

% ------------------------------------------------------------
\subsection{Visão geral das redes}\label{sec:res-geral}

\previsto{Uma tabela com as métricas exigidas pela disciplina para todas as redes da
configuração principal (vocabulário, \cod{lex}, \cod{L01}, \cod{L18}, \cod{L36}): $n$, $m$,
grau médio, densidade, reciprocidade, clusterização global e local média, distância média,
diâmetro, número de componentes e tamanho do maior, e assimetria do grau de entrada. Um
parágrafo compara as camadas: o que é esperado por construção (por exemplo, $\langle k\rangle$
entre $k$ e $2k$) e o que muda de fato. Dados: etapa \etapa{metrics}.}

\tabelaopcional[etapa \etapa{report} a partir de \etapa{metrics}]{metricas_globais}
    {Métricas estruturais das redes da configuração principal ($k=10$, união; média $\pm$
    desvio entre as 20 sementes na rede lexical, semente 0 nas camadas contextuais).}
    {tab:metricas-globais}
\figuraopcional[etapa \etapa{report} a partir de \etapa{knn} e \etapa{metrics}]{graus_camadas}
    {Distribuição do grau de entrada (CCDF, escala log-log) nas redes por ocorrência
    direcionadas, por representação. A linha vertical marca $k$, a média comum a todas.}
    {fig:graus-camadas}
\figuraopcional[etapa \etapa{report} a partir de \etapa{knn} e \etapa{metrics}]{subgrafo_camadas}
    {O mesmo subgrafo (vértices de um parágrafo do núcleo e seus vizinhos) desenhado com
    posições fixas nas redes \cod{lex}, \cod{L01}, \cod{L18} e \cod{L36}; cor pelo tipo de
    token.}
    {fig:subgrafo-camadas}

% ------------------------------------------------------------
\subsection{Rede lexical de tipos (vocabulário)}\label{sec:res-vocab}

\previsto{Descrição da geometria do espaço de entrada com os \nTiposVocab{} tipos do
vocabulário: distribuição de graus de entrada e \emph{hubs} (quais tokens são escolhidos por
muitos: tokens pouco treinados, pedaços de bytes, tokens de outras escritas?), clusterização,
distância média e diâmetro exatos, comunidades do Leiden comparadas com a classe de escrita (NMI
e ARI, com linha de base), fração de tokens que aparecem no corpus em cada comunidade e a
vizinhança das palavras-alvo no vocabulário inteiro (os 10 vizinhos de \tok{\esp banco},
\tok{\esp campo}\ldots). Dados: etapas \etapa{knn} (\cod{vocab\_k10.npz}), \etapa{metrics} e
\etapa{analyze}.}

\tabelaopcional[etapa \etapa{report} a partir de \etapa{metrics}]{vocab_comunidades}
    {Rede do vocabulário: maiores comunidades, composição por classe de escrita e fração de
    tipos presentes no corpus; NMI e ARI entre comunidades e classe de escrita.}
    {tab:vocab-comunidades}
\tabelaopcional[etapa \etapa{report} a partir de \etapa{analyze}]{vocab_vizinhos_alvos}
    {Os 10 vizinhos mais próximos de cada palavra-alvo e de controle na rede do vocabulário.}
    {tab:vocab-vizinhos-alvos}
\figuraopcional[etapa \etapa{report} a partir de \etapa{metrics}]{vocab_graus}
    {Rede do vocabulário: CCDF do grau de entrada, separada por classe de escrita.}
    {fig:vocab-graus}
\figuraopcional[etapa \etapa{report} a partir de \etapa{metrics}]{vocab_comunidades}
    {Rede do vocabulário: composição das maiores comunidades por classe de escrita.}
    {fig:vocab-comunidades}

% ------------------------------------------------------------
\subsection{P1: quanto e de que forma a organização muda}\label{sec:res-p1}

\previsto{\textbf{Escala local.} Distribuições de $J_i(\mathrm{lex},\ell)$ e de $D_i$ em cada
camada, por faixa de frequência (amostra e corpus), categoria de token e estrato, com controle
pela faixa de posição; a dominância relativa ao teto $\tilde D$; o piso de ruído por faixa ao
lado de cada distribuição. \textbf{Escala global.} CCDF do grau de entrada e assimetria por
camada, identidade dos \emph{hubs} de cada camada (que tokens, que posições) e se eles se
mantêm entre camadas, reciprocidade, clusterização e distâncias (Tabela~\ref{tab:metricas-globais}),
número e tamanho das comunidades. Resposta à H1: a dominância cai, mas continua acima do acaso
no bloco 36? Os \emph{hubs} mudam? Dados: etapas \etapa{knn}, \etapa{metrics} e
\etapa{analyze}.}

\figuraopcional[etapa \etapa{report} a partir de \etapa{analyze}]{p1_jaccard_faixas}
    {Distribuição de $J_i(\mathrm{lex},\ell)$ para $\ell\in\{1,18,36\}$ por faixa de frequência
    na amostra, com o piso de ruído $1-M^{\text{ruído}}$ de cada faixa.}
    {fig:p1-jaccard-faixas}
\figuraopcional[etapa \etapa{report} a partir de \etapa{analyze}]{p1_dominancia}
    {Dominância lexical $D$ e dominância relativa ao teto $\tilde D$ por camada, faixa de
    frequência e categoria de token ($f_t\ge3$).}
    {fig:p1-dominancia}
\tabelaopcional[etapa \etapa{report} a partir de \etapa{analyze}]{p1_resumo}
    {P1: $J$, $M$ acima do piso, $D$ e $\tilde D$ médios por camada, estrato e categoria de
    token.}
    {tab:p1-resumo}
\tabelaopcional[etapa \etapa{report} a partir de \etapa{metrics}]{p1_hubs}
    {Os 10 vértices de maior grau de entrada em cada camada: token, categoria, faixa de
    posição e grau.}
    {tab:p1-hubs}

% ------------------------------------------------------------
\subsection{P2: em que camadas a mudança se concentra}\label{sec:res-p2}

\previsto{Mudança $M$ acima do piso e $\Delta D$ em cada transição (lexical $\to$ 1, $1\to18$,
$18\to36$) e nas acumuladas, por faixa e categoria; fração de vértices cuja maior mudança cai em
cada transição; curvas $J(\ell,\ell+1)$ e $D(\ell)$ nos 36 blocos, a partir de
\cod{all\_layers.npy}. Resposta à H2: fragmentos e palavras funcionais mudam mais cedo que
palavras inteiras de conteúdo? Dados: etapas \etapa{extract} (36 camadas), \etapa{knn} e
\etapa{analyze}.}

\figuraopcional[etapa \etapa{report} a partir de \etapa{analyze}]{p2_curvas_camadas}
    {Curvas camada a camada: $J(\ell,\ell+1)$ entre blocos consecutivos e dominância
    $D(\ell)$, por categoria de token.}
    {fig:p2-curvas-camadas}
\figuraopcional[etapa \etapa{report} a partir de \etapa{analyze}]{p2_transicoes}
    {Fração de vértices cuja maior mudança acontece em cada transição, por categoria de token
    e faixa de frequência.}
    {fig:p2-transicoes}
\tabelaopcional[etapa \etapa{report} a partir de \etapa{analyze}]{p2_transicoes}
    {P2: mudança média acima do piso e $\Delta D$ por transição, por categoria de token.}
    {tab:p2-transicoes}

% ------------------------------------------------------------
\subsection{P3: comunidades e contexto}\label{sec:res-p3}

\previsto{Só no núcleo (\nVerticesNucleo{} vértices): NMI bruto, NMI menos a linha de base,
ARI e pureza das comunidades de cada camada contra tipo, tema, artigo, parágrafo, sentença,
próximo token previsto e real, e faixa de posição (controle de artefato). Tema também na rede
inteira. Exemplos qualitativos: as comunidades mais puras por parágrafo e por próximo token no
bloco 36. Estabilidade do Leiden entre as 10 execuções e, na rede lexical, concordância entre
as partições das sementes 0 e 1 do desempate. Resposta à H3: a concordância com o tipo cai e
a com o contexto e com o próximo token sobe? A faixa de posição explica as comunidades? Dados:
etapas \etapa{metrics} (comunidades) e \etapa{analyze} (concordância).}

\figuraopcional[etapa \etapa{report} a partir de \etapa{analyze}]{p3_nmi_camadas}
    {NMI menos a linha de base de permutação entre as comunidades de cada camada e cada rótulo
    (núcleo da amostra), na semente 0 do desempate. A estabilidade do Leiden (NMI médio entre
    as 10 execuções) é informada ao lado.}
    {fig:p3-nmi-camadas}
\tabelaopcional[etapa \etapa{report} a partir de \etapa{analyze}]{p3_concordancia}
    {P3: NMI bruto, linha de base, ARI e pureza por camada e rótulo (núcleo).}
    {tab:p3-concordancia}

% ------------------------------------------------------------
\subsection{P4: separação das ocorrências de um mesmo token}\label{sec:res-p4}

\previsto{Para os tipos com $f_t\ge10$, separados em alvos, controles, multitema e palavras
funcionais: número efetivo de comunidades $e^{H_t}$, NMI comunidade$\times$tema dentro do tipo,
autossimilaridade, diferença entre cosseno intra e entre temas, por camada. Comparação entre
alvos e controles (mesma alocação de temas e cotas) para testar a H4. Separação pelo sentido
anotado em \cod{senses.csv}. Redes ego de 2 ou 3 alvos com layout fixo entre camadas. Dados:
etapas \etapa{metrics} e \etapa{analyze}; anotação manual.}

\figuraopcional[etapa \etapa{report} a partir de \etapa{analyze}]{p4_separacao}
    {Separação intra-tipo por camada: autossimilaridade, diferença entre cosseno intra e entre
    temas e NMI comunidade$\times$tema, para alvos, controles e multitema.}
    {fig:p4-separacao}
\tabelaopcional[etapa \etapa{report} a partir de \etapa{analyze}]{p4_alvos}
    {P4: medidas de separação por palavra-alvo e de controle no bloco 36, com o valor na camada
    lexical como referência.}
    {tab:p4-alvos}
\figuraopcional[etapa \etapa{report} a partir de \etapa{analyze}]{p4_ego_banco}
    {Rede ego de \tok{\esp banco} nas redes \cod{lex}, \cod{L01}, \cod{L18} e \cod{L36}, com
    as ocorrências coloridas pelo tema de sentido e posições fixas entre camadas.}
    {fig:p4-ego-banco}

% ------------------------------------------------------------
\subsection{Extensão: vizinhos no vocabulário}\label{sec:res-lente}

\previsto{Se a extensão for feita (etapa \etapa{lens}): posição mediana do próprio token e do
próximo token no ranking das linhas da matriz de \emph{embeddings}, camada a camada, por
categoria de token e estrato; Jaccard entre os vizinhos no vocabulário de camadas consecutivas;
as duas checagens da Seção~\ref{sec:met-lente}. A pergunta é em que camada a ocorrência deixa de
se parecer com o próprio token e passa a se parecer com o próximo.}

\figuraopcional[etapa \etapa{report} a partir de \etapa{lens}]{lente_ranking}
    {Posição mediana do próprio token e do próximo token entre os vizinhos no vocabulário, por
    camada e categoria de token.}
    {fig:lente-ranking}

% ============================================================
\section{Robustez}\label{sec:robustez}
% ============================================================

As conclusões das Seções~\ref{sec:res-p1}--\ref{sec:res-p4} só valem se não dependerem de uma
escolha particular de construção. A grade de variações está na Tabela~\ref{tab:grade}; cada
variação muda uma dimensão por vez em relação à configuração principal. Esta seção ainda é só
estrutura.

\previsto{Uma tabela automática que, para cada variação, repete as medidas-resumo principais
(assimetria do grau de entrada, clusterização, distância média, $J$ médio acima do piso,
$D$ médio, NMI menos a linha de base contra tipo e contra parágrafo, separação dos alvos menos a
dos controles) e marca quando a conclusão qualitativa muda em relação à configuração principal.
Dados: etapas \etapa{knn} (todas as variantes), \etapa{metrics} e \etapa{analyze}.}

\tabelaopcional[etapa \etapa{report} a partir de \etapa{metrics} e \etapa{analyze}]{robustez}
    {Medidas-resumo em cada variação de robustez, ao lado da configuração principal.}
    {tab:robustez}

\subsection{Número de vizinhos e simetrização}
\previsto{$k\in\{5,20\}$ contra $k=10$, e a versão mútua contra a união: as tendências entre
camadas se mantêm? Na mútua, quantos componentes surgem e quanto a assimetria do grau cai
(a mútua corta as arestas que chegam a \emph{hubs}).}

\subsection{Tratamento de empates e tolerância}
\previsto{Variantes (b) (todos os empatados) e (c) (tipos distintos) e desempate por posição,
contra o sorteio; $\eps=10^{-4}$ contra $10^{-6}$ (arquivos \cod{nbr\_<rep>\_k10\_eps.npz}). Em
particular: quanto do $J$ lexical $\to$ bloco 1 some quando o ruído do desempate é eliminado
(variante (b), $J^w$), e quantos empates novos aparecem nas camadas contextuais com a tolerância
maior.}

\figuraopcional[etapa \etapa{report} a partir de \etapa{analyze}]{robustez_empates}
    {$J(\mathrm{lex},\ell)$ médio nas variantes de empate (a), posição, (b) e (c), e o
    Jaccard por composição $J^w$, por faixa de frequência.}
    {fig:robustez-empates}

\subsection{Normalização final e centragem}
\previsto{\cod{L36n} contra \cod{L36}: quanto a RMSNorm final (ganho por dimensão) muda as
vizinhanças. \cod{L01c}, \cod{L18c}, \cod{L36c} contra as versões sem centragem: se a
anisotropia e as ativações massivas (Seção~\ref{sec:mod-diagnosticos}) estiverem dominando o
cosseno, a centragem deve mudar sobretudo os \emph{hubs}.}

\subsection{Resolução das comunidades e subconjunto do núcleo}
\previsto{Resoluções $\gamma\in\{0{,}5;\,2\}$ do Leiden e o Louvain como checagem: a ordem dos
rótulos na P3 se mantém? Redes construídas só com o núcleo ($\approx$11 mil vértices): os
estratos esparsos alteram a estrutura global?}

\textbf{Planejado, ainda sem etapa.} As redes só do núcleo \emph{não} podem ser obtidas
restringindo as redes da amostra inteira. Pelo argumento da Seção~\ref{sec:vocab-vs-amostra}, o
$k$-NN de um subconjunto não é o subgrafo induzido do $k$-NN completo: quando vizinhos de fora
do núcleo são removidos, cada vértice precisa escolher novos vizinhos dentro dele. Restringir
os candidatos gravados (\cod{cand\_<rep>.npz}) às colunas do núcleo também não serve, porque uma
linha pode ficar com menos de $k$ candidatos e perde a garantia de completude. É preciso um
$k$-NN próprio sobre as linhas do núcleo, com candidatos recalculados; nenhuma etapa o produz
ainda (seria uma opção da etapa \etapa{knn} restrita ao estrato \cod{core}), e este item fica
de fora da tabela de robustez até lá.

% ============================================================
\section{Discussão}\label{sec:discussao}
% ============================================================

Esta seção será escrita depois dos resultados e da robustez. Ela não traz números novos: junta
o que as Seções~\ref{sec:resultados} e~\ref{sec:robustez} mostram e responde às perguntas.

\previsto{
\begin{itemize}[leftmargin=1.2em]
    \item \textbf{Resposta a cada pergunta (P1--P4)}, confrontada com a hipótese
    correspondente (H1--H4) e com o grau de confiança dado pela robustez.
    \item \textbf{O que as redes mostram além das estatísticas de pares.} Por exemplo, se a
    queda da autossimilaridade~\cite{ethayarajh2019} vem acompanhada de troca de vizinhos
    locais, de novos \emph{hubs} ou de comunidades organizadas por contexto.
    \item \textbf{Papel da fragmentação do Qwen3.} O que muda entre palavras inteiras e
    fragmentos, e o que isso implica para a generalização a outros tokenizers.
    \item \textbf{Bloco 36 e previsão do próximo token.} Se a última camada se organiza pelo
    próximo token, e como isso se relaciona com os \emph{embeddings} amarrados
    (Seção~\ref{sec:mod-escolha}) e com a extensão ``vizinhos no vocabulário''.
    \item \textbf{Rede do vocabulário.} O que a geometria do espaço de entrada inteiro diz sobre
    os tipos que aparecem no português e sobre tokens pouco treinados.
    \item \textbf{Comparação com o piloto} (GPT-2 small, 104 tokens): o que se confirmou em
    escala e o que era artefato do desempate por posição.
\end{itemize}}

% ============================================================
\section{Limitações}\label{sec:limitacoes}
% ============================================================

As limitações abaixo são as conhecidas até o planejamento e a etapa \etapa{corpus}. Cada uma diz
o que pode distorcer os resultados e o que foi feito para reduzir o efeito. A lista é
atualizada a cada etapa.

\begin{description}
    \item[O tema é só um indicador aproximado do sentido.] Na P4, o tema do artigo faz o papel
    do sentido de uma palavra-alvo: \emph{banco} em Computação é, em geral, banco de dados, mas
    um artigo de Computação pode falar do Banco Central. Uma separação por tema pode, então,
    subestimar a separação por sentido (ocorrências com o ``sentido errado'' para o tema) ou
    refletir diferenças de tema que não são de sentido (vocabulário vizinho). Mitigações: a
    anotação manual de cerca de 20 ocorrências por alvo (\cod{senses.csv}), analisada ao lado do
    tema, e a comparação com os controles, que têm a mesma alocação por tema mas um sentido
    estável: o que os controles mostrarem de ``separação por tema'' é o efeito do tema sem
    mudança de sentido.

    \item[O Qwen3 fragmenta o português.] Só 28\% dos vértices do estudo de viabilidade são
    palavras inteiras (Seção~\ref{sec:viabilidade}). Com isso, boa parte dos vértices é
    fragmento sem significado próprio, cuja vizinhança lexical reflete a ortografia e não o
    sentido, e as palavras-alvo ficaram restritas às que são um token só (12 de 16; saíram
    \emph{corrente}, \emph{onda}, \emph{célula} e \emph{núcleo}). As conclusões sobre
    fragmentos e palavras inteiras são reportadas separadamente, e as interpretações
    semânticas se restringem a palavras inteiras de conteúdo. A generalização para modelos com
    tokenizers melhores para o português (Tucano, GPT-2 português) não é testada.

    \item[Predefinições removidas deixam frases incompletas.] A primeira frase de muitos
    artigos usa predefinições para pronúncia, etimologia, datas e nomes em outras línguas
    (``X (em grego: \ldots) é\ldots''). Elas são descartadas na limpeza, e a limpeza tira os
    restos visíveis (parênteses vazios, vírgulas repetidas), mas algumas frases ficam mais
    curtas ou com construções pouco naturais. O modelo vê esse texto como contexto, então parte
    das primeiras posições dos parágrafos iniciais tem um contexto um pouco diferente do texto
    original. O efeito deve ser pequeno (só atinge algumas frases iniciais), e a faixa de
    posição entra como controle nas análises.

    \item[A busca por categorias é truncada.] A árvore de categorias da Wikipédia em português
    não é uma taxonomia limpa: em poucos níveis, uma categoria de ciência chega a pessoas,
    lugares e obras. Por isso a busca vai só até a profundidade 2 e para em 600 títulos por
    tema. O corte não é aleatório. A API lista os membros de cada categoria na ordem da chave de
    ordenação (em geral, alfabética), e a busca em largura abre as subcategorias nessa ordem,
    nível por nível; quando o limite é atingido, as subcategorias que ainda estão na fila nunca
    são abertas. O embaralhamento com semente (Seção~\ref{sec:corpus}) vem depois e só escolhe
    os 150 artigos de cada tema entre esses 600 títulos. Na execução principal, os oito temas
    atingiram o limite, e a busca abriu só uma parte das subcategorias que encontrou: de 7 de
    79 (Computação) a 98 de 439 (Música); em Política, 23 de 443. Em Biologia, Política e
    Computação, o corte veio antes de terminar a profundidade 1, e nenhum título veio da
    profundidade 2; em Computação, 494 dos 600 títulos são membros diretos da raiz. Em Música e
    Futebol, a busca parou no meio de uma lista de subcategorias por país (em ``Música da
    Austrália'' e ``Futebol das Bahamas''), de modo que, dessa lista, só os países do começo do
    alfabeto foram abertos. (Números obtidos refazendo a busca a partir do cache da etapa
    \etapa{corpus}, sem rede; a busca refeita reproduz a categoria e a profundidade de todos os
    artigos examinados.) Ficaram de fora subáreas inteiras: em Biologia, 17 das 33
    subcategorias diretas da raiz, como Reprodução, Simbiose e Teorias biológicas; em Política,
    33 de 55, como Movimentos sociais e Políticas públicas; em Computação, 43 de 49, como
    Inteligência artificial, Hardware e Mineração de dados (uma delas, Redes de computadores,
    entra como categoria extra). Entre os 150 artigos mantidos de
    Futebol, 61 vêm da profundidade 2: 32 de categorias por continente e 29 de países de
    Afeganistão a Azerbaijão; não há nenhum de Brasil ou Portugal. A profundidade dos artigos
    mantidos também varia de um tema para outro: em Computação, 123 dos 150 são membros diretos
    da raiz; em Geografia, 105 dos 150 vêm da profundidade 2. Os temas representam, portanto,
    as subcategorias mais próximas da raiz e alfabeticamente primeiras, e não a área inteira.
    Como o objetivo é ter temas bem separados, e não estimar o vocabulário de cada área, isso é
    aceitável, mas os nomes dos temas devem ser lidos como ``artigos próximos da categoria
    raiz''. A comparação entre representações não é afetada, porque os vértices são as mesmas
    ocorrências em todas as redes; as análises que usam o tema como rótulo (P3, e P4 com o tema
    como indicador do sentido) herdam o recorte. O mesmo corte limita o
    descarte de páginas multitema: um título só é reconhecido como multitema se estiver na
    lista truncada de mais de um tema, e um artigo que também pertence à parte não aberta de
    outra árvore passa como se fosse de um tema só. O descarte de páginas multitema e de
    biografias reduz a mistura, ao custo de excluir páginas de fronteira entre temas.

    \item[Precisão \emph{bfloat16}.] O modelo roda em bf16 (8 bits de mantissa), e os vetores
    são guardados em bf16. O cosseno em \emph{float64} é exato \emph{para os vetores
    guardados}, mas esses vetores já carregam o arredondamento das contas do modelo: duas
    execuções com sequências de comprimentos diferentes podem dar estados ligeiramente
    diferentes para o mesmo prefixo (Seção~\ref{sec:mod-sequencias}). Vizinhanças decididas por
    diferenças de cosseno menores que esse ruído são, na prática, empates. Mitigações: a cópia
    por grupo de prefixo (que torna exatos os empates que deveriam ser exatos), o diagnóstico da
    maior diferença dentro dos grupos (\difMaxPrefixo) e o teste de sensibilidade com
    $\eps=10^{-4}$. Na GPU da execução (RTX A5500, 24 GiB), o modelo caberia em
    \emph{float32}, mas isso não foi adotado: o bf16 é o tipo nativo do modelo, os vetores
    continuam guardados em bf16, e o passo de arredondamento do bf16 nas magnitudes do bloco 36
    já é da ordem da diferença medida dentro dos grupos de prefixo.

    \item[A Wikipédia provavelmente está no pré-treino.] O Qwen3 foi pré-treinado em cerca de 36
    trilhões de tokens de fontes não detalhadas~\cite{qwen3}, e a Wikipédia em português quase
    certamente faz parte delas. O modelo pode, então, ter visto (e em parte memorizado) os
    parágrafos analisados. Isso não invalida a análise da geometria das representações, mas
    pode torná-las mais ``confiantes'' do que seriam em texto novo (por exemplo, uma previsão
    do próximo token mais acertada e, no bloco 36, uma organização mais forte pelo próximo
    token). Não há como controlar isso sem texto certamente fora do pré-treino; fica registrado
    como ressalva de interpretação.

    \item[Um modelo, uma língua, uma amostra.] Os resultados são de um único modelo
    (Qwen3-4B-Base, numa revisão fixada), de texto enciclopédico em português e de uma amostra
    desenhada (as frequências na amostra refletem as cotas, e não o texto natural). As faixas
    de frequência no corpus e a análise da P3 só no núcleo reduzem o efeito do desenho, mas
    não o eliminam.
\end{description}

% ============================================================
\section{Problemas encontrados e soluções}\label{sec:problemas}
% ============================================================

Cada problema é registrado com o \textbf{sintoma} (o que se observou), a \textbf{causa}, a
\textbf{solução} adotada e o \textbf{impacto} nos resultados. O objetivo é que quem refizer o
experimento não tropece nos mesmos pontos e que fique claro quais escolhas do método nasceram de
um problema concreto. As entradas abaixo são da fase de planejamento e das primeiras etapas.

\newcommand{\problema}[5]{%
    \subsection{#1}\label{#2}
    \begin{description}[leftmargin=1.2em, itemsep=1pt]
        \item[Sintoma.] #3
        #4
        \item[Impacto.] #5
    \end{description}}

\problema{Limitação de taxa da API da Wikipédia (HTTP 429)}{sec:prob-429}
{No estudo de viabilidade, depois de algumas centenas de requisições seguidas, a API passou a
responder HTTP 429 (\emph{Too Many Requests}) e o download parava no meio.}
{\item[Causa.] A política da Wikimedia limita a taxa de requisições por cliente, sobretudo de
clientes sem um \cod{User-Agent} identificado; buscar artigos um a um multiplica o número de
requisições.
\item[Solução.] O cliente da etapa \etapa{corpus} (\cod{wiki.py}) pede o \emph{wikitext} em
lotes de 50 páginas por requisição (o máximo da API), espera 2,5\,s entre requisições, envia um
\cod{User-Agent} com contato (configurável), usa \cod{maxlag=5}, respeita o cabeçalho
\cod{Retry-After} e, na falta dele, recua exponencialmente (30\,s, 60\,s, \ldots, até 10
minutos). Toda resposta crua fica num cache em disco, então um download interrompido recomeça
de onde parou e a limpeza pode ser refeita sem rede. Os artigos são fixados por \cod{revid} no
\cod{manifest.csv}.}
{Nenhum nos resultados. O download completo fica mais lento (horas, em segundo plano), mas é
retomável e reprodutível.}

\problema{Páginas de usuário e categorias de manutenção na busca por categorias}
{sec:prob-categorias}
{As primeiras buscas nas árvores de categorias traziam páginas de usuário, predefinições,
listas e páginas como ``Artigos sem fontes'' e ``Esboços sobre física''.}
{\item[Causa.] A árvore de categorias da Wikipédia mistura categorias de conteúdo com
categorias de manutenção (artigos que precisam de revisão, esboços, páginas de discussão) e de
organização interna, que também são subcategorias das raízes temáticas.
\item[Solução.] Só entram membros do domínio principal (artigos, \cod{ns=0}); subcategorias cujo
título contém ``Artigos'', ``Esboço'', ``Páginas'', ``Predefinições'', ``Wikipédia'',
``Usuário'', ``Listas'' ou ``Categoria:!'' não são seguidas; títulos ``Lista de\ldots'' são
descartados; a profundidade é limitada a 2 (Seção~\ref{sec:corpus}).}
{Os temas ficam mais limpos. A lista de padrões é heurística e pode deixar passar alguma
categoria de manutenção com outro nome; o descarte de páginas multitema e de biografias
reduz o que sobra.}

\problema{Fragmentação do português no tokenizer do Qwen3}{sec:prob-fragmentacao}
{No estudo de viabilidade, o Qwen3 produziu 1,78 tokens por palavra, e só 28\% dos vértices
eram palavras inteiras (54\% no Tucano, 72\% no GPT-2 português). Quatro das 16 palavras-alvo
candidatas (\emph{corrente}, \emph{onda}, \emph{célula}, \emph{núcleo}) não são um token só.}
{\item[Causa.] O vocabulário BPE do Qwen3 (151.669 ids) é multilíngue e dominado por inglês,
chinês e código; palavras portuguesas menos frequentes, principalmente com acentos, são
divididas em pedaços.
\item[Solução.] As palavras-alvo ficaram restritas às que são um token só (12), contadas só no id
minúsculo precedido de espaço. As categorias de token (Seção~\ref{sec:categorias}) entram em
todas as análises, e as interpretações semânticas se concentram em palavras inteiras de
conteúdo. O custo foi pesado na escolha do modelo (Seção~\ref{sec:mod-escolha}).}
{Menos palavras-alvo e uma fração grande de vértices que são fragmentos. É também um resultado
em si: a P2 compara fragmentos e palavras inteiras.}

\problema{\texttt{hidden\_states[36]} normalizado e \texttt{torch\_dtype} obsoleto no
transformers 5}{sec:prob-hidden}
{Ao planejar a extração, a documentação e o código do transformers 5.16 mostraram que o último
elemento de \cod{hidden\_states} não é a saída bruta do bloco 36. Além disso, carregar o modelo
com \cod{torch\_dtype=} emite um aviso de argumento obsoleto.}
{\item[Causa.] Nos modelos decodificadores do transformers 5, a tupla \cod{hidden\_states}
recebe como último elemento o estado \emph{depois} da RMSNorm final (o que a cabeça de
linguagem lê), no lugar da saída bruta do último bloco; os demais elementos são saídas brutas.
O argumento de tipo passou a se chamar \cod{dtype}.
\item[Solução.] As saídas brutas são capturadas por \emph{hooks} em \cod{layers[0]},
\cod{layers[17]} e \cod{layers[35]} (e nos 36 blocos para o \emph{memmap}), e a versão
normalizada por um \emph{hook} em \cod{norm}, que vira a variante de robustez \cod{L36n}. A
verificação \cod{--verify} confere \cod{norm(hook\_36) == hidden\_states[36]} no modelo real
(Seção~\ref{sec:mod-verificacao}). O carregamento usa \cod{dtype=torch.bfloat16}.}
{Sem os \emph{hooks}, o bloco 36 seria comparado aos blocos 1 e 18 num espaço diferente
(normalizado e reescalado dimensão a dimensão), e a transição $18\to36$ misturaria o efeito do
bloco com o da normalização.}

\problema{Limite por ordem de chegada e artigo igual a parágrafo na simulação}
{sec:prob-simulacao}
{Na simulação do desenho híbrido (\cod{sim\_hybrid.py}), as ocorrências mantidas dos tipos
frequentes (``de'', ``,'') vinham todas dos primeiros parágrafos do núcleo, e cada parágrafo do
núcleo vinha de um artigo diferente.}
{\item[Causa.] O limite $f_{\max}$ era aplicado durante a varredura: a partir da 50.ª ocorrência
de um tipo, as seguintes eram ignoradas. E o núcleo escolhia um parágrafo por artigo.
\item[Solução.] No desenho final (Seção~\ref{sec:amostra}), o núcleo é montado primeiro e o
limite é aplicado depois, por sorteio uniforme entre todas as ocorrências de cada tipo acima do
limite. O núcleo usa 5 artigos $\times$ 2 parágrafos por tema. Um teste verifica que o limite é
aplicado por sorteio.}
{Com o limite por ordem de chegada, os tipos frequentes ficariam concentrados em poucos
parágrafos, o que enviesaria a P3 (tipo e parágrafo ficariam confundidos). Com um parágrafo por
artigo, os rótulos ``artigo'' e ``parágrafo'' da P3 seriam praticamente iguais e não poderiam ser
distinguidos.}

\problema{Conteúdo de referências no texto limpo (\texttt{mwparserfromhell} 0.7.2)}
{sec:prob-ref}
{No corpus do estudo de viabilidade, limpo só com \cod{strip\_code}, apareciam no meio dos
parágrafos títulos de livros, nomes de autores, URLs e datas de acesso, colados à palavra
anterior (``A energiaSilva, Livro, 2001 é\ldots''), e frases com buracos no lugar de fórmulas.}
{\item[Causa.] Na versão 0.7.2 do \cod{mwparserfromhell}, o \cod{strip\_code} remove a marca
\cod{<ref>}, mas mantém o \emph{conteúdo} dela como texto; já o conteúdo de \cod{<math>} some
sem deixar nada no lugar (``a energia é dada por , onde\ldots'').
\item[Solução.] Antes da conversão, a árvore sintática é reescrita: marcas \cod{<ref>} e outras
de bloco são removidas com o conteúdo, e linhas com \cod{<math>} e afins são descartadas
(Seção~\ref{sec:corpus}). Testes com trechos reais de \emph{wikitext} cobrem esses casos.}
{Os números do estudo de viabilidade vêm do corpus com esse defeito; são aproximados, mas o
defeito não muda as conclusões do estudo (ordens de grandeza de parágrafos, tokens e
repetições). O corpus final não tem o problema.}

\problema{\emph{Virtualenv} copiado com o \texttt{pytest} de outro projeto}{sec:prob-venv}
{\cod{.venv/bin/pytest} falhava com erros de importação ou rodava com o interpretador errado,
embora \cod{.venv/bin/python -m pytest} funcionasse.}
{\item[Causa.] O ambiente virtual foi copiado de outro projeto. Os \emph{scripts} de entrada
(\cod{pytest}, \cod{ruff}\ldots) têm na primeira linha (\emph{shebang}) o caminho absoluto do
interpretador, que continuava apontando para o \cod{.venv} do projeto original.
\item[Solução.] Rodar as ferramentas pelo interpretador do projeto
(\cod{.venv/bin/python -m pytest}) ou pelo \cod{uv run}, que usa o ambiente certo; recriar o
ambiente com \cod{uv sync --dev} resolve de vez.}
{Nenhum nos resultados. Fica o alerta de sempre conferir qual interpretador executa os testes.}


% ------------------------------------------------------------
\begin{thebibliography}{99}
\small

\bibitem{blondel2008}
V.~D. Blondel, J.-L. Guillaume, R.~Lambiotte e E.~Lefebvre.
Fast unfolding of communities in large networks.
\emph{Journal of Statistical Mechanics: Theory and Experiment}, P10008, 2008.

\bibitem{csardi2006}
G.~Csárdi e T.~Nepusz.
The igraph software package for complex network research.
\emph{InterJournal, Complex Systems}, 1695, 2006.

\bibitem{ethayarajh2019}
K.~Ethayarajh.
How contextual are contextualized word representations? Comparing the geometry of BERT,
ELMo, and GPT-2 embeddings.
Em \emph{Proceedings of EMNLP-IJCNLP}, 2019.

\bibitem{hubert1985}
L.~Hubert e P.~Arabie.
Comparing partitions.
\emph{Journal of Classification}, 2(1):193--218, 1985.

\bibitem{newman2006}
M.~E.~J. Newman.
Modularity and community structure in networks.
\emph{Proceedings of the National Academy of Sciences}, 103(23):8577--8582, 2006.

\bibitem{press2017}
O.~Press e L.~Wolf.
Using the output embedding to improve language models.
Em \emph{Proceedings of EACL}, 2017.

\bibitem{qwen3}
Qwen Team.
Qwen3 technical report.
\emph{arXiv:2505.09388}, 2025.

\bibitem{radovanovic2010}
M.~Radovanović, A.~Nanopoulos e M.~Ivanović.
Hubs in space: popular nearest neighbors in high-dimensional data.
\emph{Journal of Machine Learning Research}, 11:2487--2531, 2010.

\bibitem{sun2024}
M.~Sun, X.~Chen, J.~Z. Kolter e Z.~Liu.
Massive activations in large language models.
\emph{arXiv:2402.17762}, 2024.

\bibitem{timkey2021}
W.~Timkey e M.~van Schijndel.
All bark and no bite: rogue dimensions in transformer language models obscure
representational quality.
Em \emph{Proceedings of EMNLP}, 2021.

\bibitem{traag2019}
V.~A. Traag, L.~Waltman e N.~J. van Eck.
From Louvain to Leiden: guaranteeing well-connected communities.
\emph{Scientific Reports}, 9:5233, 2019.

\bibitem{vinh2010}
N.~X. Vinh, J.~Epps e J.~Bailey.
Information theoretic measures for clusterings comparison: variants, properties,
normalization and correction for chance.
\emph{Journal of Machine Learning Research}, 11:2837--2854, 2010.

\bibitem{zhang2019}
B.~Zhang e R.~Sennrich.
Root mean square layer normalization.
Em \emph{Advances in Neural Information Processing Systems 32}, 2019.

\end{thebibliography}

\appendix
% ============================================================
% Apêndices (incluídos depois de \appendix)
% ============================================================

\section{Registro de decisões}\label{sec:decisoes}

Cada decisão de método entra aqui com o contexto em que foi tomada, as alternativas
consideradas, a escolha, o motivo e a consequência para o experimento. O registro começa com as
decisões do planejamento e cresce a cada etapa. As decisões mais finas (limiares, listas de
categorias) ficam nos arquivos de configuração e nas seções correspondentes.

{\footnotesize
\setlength{\tabcolsep}{3.5pt}
\renewcommand{\arraystretch}{1.15}
\begin{longtable}{@{}>{\raggedright\arraybackslash}p{0.035\linewidth}
    >{\raggedright\arraybackslash}p{0.155\linewidth}
    >{\raggedright\arraybackslash}p{0.155\linewidth}
    >{\raggedright\arraybackslash}p{0.17\linewidth}
    >{\raggedright\arraybackslash}p{0.21\linewidth}
    >{\raggedright\arraybackslash}p{0.185\linewidth}@{}}
\caption{Registro de decisões.}\label{tab:decisoes}\\
\toprule
\textbf{\#} & \textbf{Contexto} & \textbf{Alternativas} & \textbf{Escolha} & \textbf{Motivo} &
\textbf{Consequência} \\
\midrule
\endfirsthead
\multicolumn{6}{@{}l}{\emph{Tabela~\ref{tab:decisoes}, continuação}}\\
\toprule
\textbf{\#} & \textbf{Contexto} & \textbf{Alternativas} & \textbf{Escolha} & \textbf{Motivo} &
\textbf{Consequência} \\
\midrule
\endhead
\midrule
\multicolumn{6}{r@{}}{\emph{continua}}\\
\endfoot
\bottomrule
\endlastfoot
D1 & Modelo a analisar, entre 4 e 8 bilhões de parâmetros e cabendo numa GPU de 8 GiB. &
Qwen3-4B pós-treinado; Tucano-2b4; GPT-2 small português (piloto). &
Qwen3-4B-Base, revisão fixada por SHA. &
É o modelo de linguagem puro (sem ajuste para diálogo), com arquitetura atual e 36 blocos para a
curva camada a camada. O Tucano tokeniza melhor, mas é menos representativo; o GPT-2 é pequeno e
foi treinado no mesmo texto. &
Português fragmentado (28\% de palavras inteiras): alvos restritos a um token e medidas sempre
por categoria (Seções~\ref{sec:mod-escolha} e~\ref{sec:prob-fragmentacao}). \\
\midrule
D2 & Amostrar cerca de 15 mil vértices que sirvam a P1--P4. &
Parágrafos inteiros (amostragem natural); natural começando pelos alvos. &
Desenho híbrido: núcleo de parágrafos inteiros e estratos de alvos, controles e multitema. &
Parágrafos inteiros deixam só 4 a 9 palavras de conteúdo repetidas em 3 temas e nenhum alvo
utilizável; o híbrido dá cerca de 170 (Seção~\ref{sec:viabilidade}). &
Frequências na amostra refletem o desenho (faixas também pelo corpus); P3 só no núcleo; cerca
de 330 mil tokens a processar. \\
\midrule
D3 & Que vértices tem a rede lexical de tipos. &
Só os $\approx$3,3 mil tipos da amostra; o vocabulário inteiro. &
Vocabulário inteiro (151.643 tipos, sem especiais nem enchimento). &
Descreve a geometria do espaço de entrada sem depender da amostra, e é uma rede grande com
estrutura própria para as métricas da disciplina. &
Uma tabela auxiliar de tipos da amostra continua necessária (não é um recorte;
Seção~\ref{sec:vocab-vs-amostra}); distâncias exatas só na configuração principal. \\
\midrule
D4 & Precisão do cosseno com tolerância de empate $\eps=10^{-6}$ e $d=2560$. &
\emph{float32} na GPU; \emph{float64} na CPU. &
\emph{float64} na CPU, em blocos. &
Em \emph{float32} o erro é de $10^{-6}$ a $10^{-5}$, comparável a $\eps$; empates falsos e
desfeitos mudariam as vizinhanças. &
Cálculo mais lento e cerca de 3 GB de memória na rede do vocabulário; empates exatos e
reprodutíveis. \\
\midrule
D5 & Qual versão do bloco 36 usar, dado que \cod{hidden\_\allowbreak states[36]} é
normalizado no transformers 5.16. &
\cod{hidden\_\allowbreak states[36]} (pós-RMSNorm); saída bruta do bloco. &
Saída bruta por \emph{hook}; a normalizada como robustez (\cod{L36n}). &
Os blocos 1 e 18 só existem na forma bruta; comparar objetos do mesmo tipo na transição
$18\to36$. &
Verificação \cod{norm(hook\_36) == hidden\_states[36]} no modelo real; uma variante a mais na
robustez (Seção~\ref{sec:prob-hidden}). \\
\midrule
D6 & Desempate dos vizinhos na fronteira do $k$-ésimo. &
Só posição (piloto); sorteio; todos os empatados; tipos distintos. &
Sorteio com 20 sementes como principal; as outras três como referência e robustez. &
A posição cria \emph{hubs} artificiais a 15 mil vértices (todas as ocorrências de um tipo
escolhem os mesmos vizinhos). &
Métricas estruturais da rede lexical com média e desvio entre sementes (semente 0 nas
contextuais, onde os empates são raros); piso de ruído e $J^w$ para separar o efeito do
sorteio (Seção~\ref{sec:red-empates}). \\
\midrule
D7 & Biblioteca de redes e detecção de comunidades. &
NetworkX (piloto); igraph; Louvain; Leiden com CPM. &
igraph, com Leiden maximizando a modularidade, 10 sementes; Louvain como checagem. &
Distâncias exatas em C (segundos por rede, contra 20 a 25 minutos no NetworkX); o Leiden
garante comunidades conexas; o padrão CPM do igraph tem outra resolução, por isso o código
passa \cod{objective\_function=\allowbreak"modularity"}. &
NetworkX fica como referência nos testes; estabilidade das comunidades reportada. \\
\midrule
D8 & O que fazer com o piloto e com o código contrafactual de gênero herdado. &
Estender o piloto; reescrever tudo; congelar. &
Piloto congelado (rodando e com testes); código contrafactual e notebook removidos (ficam no
histórico do \cod{git}). &
O piloto é monolítico e não escala, mas documenta os achados que motivaram o método; o código
contrafactual não serve às perguntas atuais. &
Pipeline novo em módulos próprios; um teste de regressão exige que o k-NN novo (posição,
união) reproduza o do piloto. \\
\midrule
D9 & Como configurar execuções (principal e ensaio). &
Argumentos de linha de comando; vários YAMLs por etapa; um YAML por execução. &
Um YAML por execução (\cod{configs/\allowbreak experiment.yaml}; \cod{configs/\allowbreak mini.yaml} herda dele). &
Uma execução fica descrita por um arquivo só, gravado no \cod{\_manifest.json} de cada etapa. &
Artefatos separados por execução (\cod{outputs/\allowbreak experiment/\allowbreak <nome>/}); o ensaio não sobrescreve a
execução principal. \\
\midrule
D10 & Como documentar o experimento. &
Escrever o relatório no fim; relatório vivo. &
Relatório técnico vivo, com figuras, tabelas e números gerados pelo pipeline. &
O texto nunca diverge da última execução, e as decisões e problemas são registrados enquanto
ainda estão frescos. &
Regra de trabalho: uma etapa só está pronta com a seção do relatório atualizada e o documento
compilando. \\
\end{longtable}
}

% ------------------------------------------------------------
\section{Reprodutibilidade}\label{sec:reprodutibilidade}

\paragraph{Ambiente.} Python 3.12 com dependências fixadas em \cod{uv.lock}; instalação com
\cod{uv sync --dev}. As versões efetivamente usadas em cada etapa (Python, numpy, torch, transformers,
igraph, networkx, pandas) são gravadas no \cod{\_manifest.json} da etapa e
reunidas na Tabela~\ref{tab:versoes}. A versão do \cod{mwparserfromhell}, que afeta a
limpeza (Seção~\ref{sec:prob-ref}), ainda não é gravada nos manifestos; até lá, ela é a fixada
em \cod{uv.lock} (0.7.2). O relatório compilado indica a data e a revisão do
\cod{git} da execução que gerou os números (\dataExecucao, \revisaoGit).

\paragraph{Modelo.} \cod{Qwen/Qwen3-4B-Base}, revisão
\cod{\seqsplit{906bfd4b4dc7f14ee4320094d8b41684abff8539}} do Hugging Face (código do próprio transformers, sem código remoto).
Pesos em bf16, \cod{device\_map="auto"} e \cod{max\_memory} de 6 GiB na GPU e 10 GiB na CPU
(\emph{offload} de parte dos blocos). O \emph{offload} muda o tempo, não as contas.

\paragraph{Dados.} A lista de artigos, com \cod{pageid} e \cod{revid}, fica versionada em
\cod{data/corpus/manifest.csv}: o mesmo texto pode ser baixado de novo pelas revisões, mesmo
que os artigos tenham mudado. As respostas cruas da API ficam em cache em \cod{data/raw/wiki}.

\paragraph{Sementes.} Todas vêm do YAML da execução:
\begin{itemize}
    \item \cod{corpus.seed} $=438$: ordem de download dos artigos de cada tema;
    \item \cod{sample.seed} $=438$: um único gerador sorteia, em ordem, o núcleo, o limite por
    tipo, alvos, controles e multitema;
    \item \cod{networks.base\_seed} $=438$: a semente $r$ do desempate da representação
    \cod{rep} com $k$ vizinhos usa \cod{numpy.random.default\_rng([438, crc32(rep), k, r])},
    $r=0,\ldots,19$;
    \item Leiden: 10 execuções com sementes fixas do gerador \cod{random} do Python (usado pelo
    igraph); permutações da linha de base com gerador semeado.
\end{itemize}

\paragraph{Comandos.} Cada etapa lê os artefatos da anterior e pode ser refeita sozinha; uma
etapa cujas saídas já existem é pulada, a menos que se passe \cod{--force}.

\begin{center}
\small
\begin{tabularx}{\linewidth}{@{}>{\ttfamily\raggedright\arraybackslash}p{0.47\linewidth} X@{}}
\toprule
\normalfont\textbf{Comando} & \textbf{O que faz} \\
\midrule
uv run gender-networks corpus & baixa e limpa os artigos (\cod{data/corpus/}) \\
uv run gender-networks sample & monta a amostra híbrida (\cod{sample/}) \\
uv run gender-networks extract -{}-verify & checa as premissas no modelo real (8 sequências) \\
uv run gender-networks extract & extrai as representações (\cod{reps/}) \\
uv run gender-networks knn & candidatos e vizinhos k-NN de todas as variantes (\cod{knn/}) \\
uv run gender-networks metrics & métricas estruturais e comunidades (\cod{metrics/}) \\
uv run gender-networks analyze & análises de P1--P4 (\cod{analysis/}) \\
uv run gender-networks lens & extensão ``vizinhos no vocabulário'' (\cod{lens/}) \\
uv run gender-networks report & figuras, tabelas e \cod{numeros.tex} deste relatório \\
uv run gender-networks all & todas as etapas em ordem (sem a extensão) \\
\bottomrule
\end{tabularx}
\end{center}
Todos aceitam \cod{--config} (padrão \cod{configs/experiment.yaml}); o ensaio de ponta a ponta
usa \cod{--config configs/mini.yaml} (2 temas, cerca de 1.500 vértices), cujos artefatos e
relatório gerado ficam em \cod{outputs/experiment/mini/}. O relatório é compilado com
\cod{latexmk -pdf relatorio.tex} em \cod{report/relatorio/}. O piloto congelado roda com
\cod{uv run python scripts/run\_network\_pilot.py}.

\tabelaopcional[etapa \etapa{report} a partir dos \cod{\_manifest.json} de todas as etapas]
    {versoes}
    {Versões do software, revisão do modelo e tempos de cada etapa na última execução.}
    {tab:versoes}


\end{document}
